% Add the following subsections
% - Problem -> Motivation
% - Solution -> Summary of Solutions
% - Detailed description of the solution
% - Architecture and Design
% - Examples
% - Comparison with the other solutions, state of the art

\input{papers/template/style.tex}

\graphicspath{{papers/Causify_AIgentic_Development_System/}}

\begin{document}
  \title{Causify AIgentic Development System}
  \maketitle

  % ############################################################################
  \section{Introduction}
  Software development workflows are becoming more complex as they adapt
  to the demands of large-scale systems and modern collaborative
  development practices. As teams and codebases grow, companies face the
  challenge of organizing both effectively. When it comes to structuring
  the codebase, two main approaches emerge: monorepos and multi-repos.
  Monorepos consolidate all code into a single repository, simplifying version
  control but carrying a risk of scalability and maintainability issues.
  Conversely, multi-repos store the code in logically separated repositories,
  easier to manage and deploy but more difficult to keep in sync.

  In this paper, we propose Causify dev system, an alternative hybrid
  solution: a modular system architecture built around runnable directories.
  Although independent, these directories maintain cohesion through shared
  tooling and environments, offering a straightforward and scalable way
  to organize the codebase while ensuring reliability in development, testing,
  and deployment.

  % ############################################################################
  \section{Contributions}

  This paper presents the Causify development system, a novel approach
  to software development infrastructure that addresses the limitations
  of traditional monorepo and multi-repo architectures, development vs deployment,

  The main contributions are:

  \begin{itemize}
    \item \textbf{Human / AI development system}:
      A development system designed to enable AI agentic coding, to increase
      the efficiency of human developers, and to allow AI coding agents and
      humans to collaborate.

    \item \textbf{Runnable directories architecture}: A hybrid organizational
      model that combines the modularity of multi-repos with the environment
      consistency of monorepos, enabling independent development
      lifecycles while maintaining unified tooling and testing
      capabilities.

    \item \textbf{Containerized development workflows}: 
      A Docker-based approach to development that ensures reproducible
      environments across development, testing, and production stages.

    \item \textbf{Dockerized Executables}: A containerization pattern for
      individual development utilities that packages tools with their
      dependencies in slim Docker images, enabling execution without local
      installation.
  \end{itemize}

  This paper is organized as in the following
  TODO(gp)

  % ############################################################################
  % ##############################################################################
\section{Code Quality}

% Related documentation files:
% - helpers_root/docs/tools/linter/all.developing_linter.how_to_guide.md
%   How to develop custom linters
% - helpers_root/docs/tools/linter/all.amp_fix_md_links.explanation.md
%   Example linter for markdown links
% - helpers_root/docs/tools/all.precommit.how_to_guide.md
%   Pre-commit hooks setup and usage
% - helpers_root/docs/tools/all.code_coverage.how_to_guide.md
%   Code coverage measurement
% - helpers_root/docs/tools/all.codecov.how_to_guide.md
%   Codecov integration for coverage tracking
% - helpers_root/docs/tools/all.code_coverage_subprocess.how_to_guide.md
%   Coverage in subprocess testing
% - helpers_root/docs/tools/all.import_check.reference.md
%   Import cycle detection tools
% - helpers_root/docs/tools/git/all.git_hooks.how_to_guide.md
%   Git hooks for quality enforcement
% - helpers_root/docs/tools/git/all.git_hooks.reference.md
%   Complete git hooks reference
% - helpers_root/docs/tools/git/all.gitleaks.how_to_guide.md
%   Secret detection with gitleaks
% - helpers_root/docs/tools/git/all.gitleaks.reference.md
%   Gitleaks configuration reference
% - helpers_root/docs/code_guidelines/all.coding_style.how_to_guide.md
%   Coding style standards
% - helpers_root/docs/code_guidelines/all.coding_style_guidelines.reference.md
%   Complete style reference
% - helpers_root/docs/code_guidelines/all.code_review.how_to_guide.md
%   Code review process and quality gates
% - helpers_root/docs/code_guidelines/all.submit_code_for_review.how_to_guide.md
%   Preparing code for review
% - helpers_root/docs/code_guidelines/all.type_hints.how_to_guide.md
%   Type hints usage for better quality
% - helpers_root/docs/build_system/all.linter_gh_workflow.explanation.md
%   Automated linting in CI/CD
% - helpers_root/docs/build_system/all.semgrep_workflow.explanation.md
%   Static analysis with Semgrep

% ==============================================================================
\subsection{Motivation}
Code quality is of paramount importance for building maintainable software
systems. As codebases grow in size and complexity, maintaining consistent
standards, detecting architectural issues, and ensuring comprehensive test
coverage become increasingly challenging. Manual code reviews alone cannot
scale to catch all violations, architectural degradation, or quality
regressions before they enter the codebase. To address these challenges, we
developed an integrated code quality system that enforces standards
automatically, provides immediate feedback to developers, and maintains
visibility into code health through comprehensive measurement and reporting.

% ==============================================================================
\subsection{Summary}

% TODO(gp_ai3): The text should correspond to the topics here and convert to latex

Our integrated code quality system combines multiple enforcement layers to
maintain high standards throughout the development lifecycle.

- The pluggable linter framework provides both automatic fixes and quality
  gates, ensuring consistent formatting and catching common errors.

- Import cycle detection prevents architectural degradation by enforcing
  acyclic dependencies.

- Pre-commit hooks provide immediate feedback at commit time, catching issues
  before they enter the repository.

- Comprehensive coverage tracking with structured test suites and
  enforced thresholds ensures all code paths remain tested as the codebase
  evolves.

- This multi-layered
approach creates a robust quality enforcement system that scales with team
size and codebase complexity. By automating enforcement and providing
immediate feedback, we reduce the burden on code reviewers while maintaining
consistent standards across all contributions. The system is designed to
be transparent, providing clear diagnostics when checks fail and maintaining
audit trails of all validations. Together, these components enable teams to
move quickly while maintaining high standards.

% Make it easy to create packages
% Ensure code has same formatting (black, ruff, ...)
% Catch errors statically using linting
% Ensure code follows the same higher level rules (e.g., using AI)
% - interface style (input, output, inout)
% - testing
% Keep code in order of dependencies (interfaces first and then the rest in
% order to make reading it easy for a human)
% Docstring formatting
% Convention for branch naming to make it easy to refer to bug tracking context
% Documentation convention (diataxis)

% ==============================================================================
\subsection{The Linter Framework}
At the core of our code quality enforcement is a pluggable linting framework
that checks and modifies code to improve readability, detect bugs, and enforce
coding standards. Unlike monolithic linters, our system is built as a
collection of focused, composable linting scripts that can be executed
independently or as an integrated pipeline.

% ------------------------------------------------------------------------------
\subsubsection{Architecture}
Each linter is implemented as a Python class that inherits from
\texttt{linters.action.Action}
and implements a standardized interface:
\begin{itemize}
  \item \textbf{\texttt{\_execute()}}: Processes a file and returns a list
    of detected violations with descriptions

  \item \textbf{\texttt{check\_if\_possible()}}: Determines if the linter
    can run in the current environment

  \item \textbf{\texttt{run()}}: Orchestrates execution across multiple files
\end{itemize}

Linters are categorized into two types based on their behavior:
\begin{itemize}
  \item \textbf{Modifying actions}: Automatically fix issues by rewriting
    files (e.g., code formatters, import sorters, whitespace normalizers)

  \item \textbf{Non-modifying actions}: Report violations without changing
    files (e.g., static analysis tools, style checkers, complexity
    analyzers)
\end{itemize}

This separation enables different execution strategies: modifying linters
run during development to automatically maintain standards, while non-modifying
linters act as quality gates in CI/CD pipelines.

% ------------------------------------------------------------------------------
\subsubsection{Comprehensive Linting Rules}

Our linting framework integrates both third-party tools and custom rules
tailored to our coding conventions. The modifying linters include:

\begin{itemize}
  \item \textbf{Code formatting}: ty, Black for consistent Python formatting,
    Prettier for Markdown files

  \item \textbf{Import management}: isort for import ordering, autoflake
    for removing unused imports

  \item \textbf{Structural consistency}: Class method ordering,
    separating line normalization, whitespace standardization

  \item \textbf{Documentation}: Docstring formatting, table of contents
    generation for notebooks
\end{itemize}

Non-modifying linters enforce additional constraints:

\begin{itemize}
  \item \textbf{Static analysis}: flake8, pylint, and mypy for catching bugs
    and type errors

  \item \textbf{Policy enforcement}: File size limits, naming conventions,
    merge conflict detection

  \item \textbf{Convention checking}: TODO format validation, import
    pattern verification, shebang requirements
\end{itemize}

% ------------------------------------------------------------------------------
\subsubsection{Integration and Extensibility}
Adding a new linting rule requires creating a Python script in the
\texttt{linters/} directory, registering it in \texttt{linters/base.py}, and
adding unit tests. This modular design allows teams to customize enforcement
without modifying core infrastructure. Linters can be executed selectively
based on file types, directories, or specific development phases, providing
flexibility while maintaining consistency.

% ==============================================================================
\subsection{Import Cycle Detection}
Circular dependencies between modules represent a significant
architectural anti-pattern that degrades code maintainability, complicates
testing, and can lead to subtle runtime errors. As codebases grow organically,
developers may inadvertently introduce import cycles that violate
architectural boundaries. To maintain clean architecture, we developed
automated tools for detecting and visualizing circular dependencies.

% ------------------------------------------------------------------------------
\subsubsection{Dependency Analysis Tools}
Our system provides two complementary tools for import analysis:
\begin{itemize}
  \item \textbf{show\_imports}: A visualization tool that generates dependency
    graphs showing relationships between modules and packages. It
    supports multiple visualization modes including file-level
    dependencies, directory-level aggregation, external package
    dependencies, and cycle-only views. The tool can limit analysis
    depth, filter by scope, and export results in various image formats.

  \item \textbf{detect\_import\_cycles}: An enforcement tool that analyzes
    the dependency graph and fails if circular dependencies are detected.
    When cycles are found, it logs all participating files in each cycle,
    enabling developers to quickly identify and resolve architectural
    violations.
\end{itemize}

Both tools are built on Python's import machinery and require that all
directories containing Python files are proper modules (containing
\texttt{\_\_init\_\_.py}).
This constraint ensures comprehensive analysis and prevents silent
failures from malformed package structures.

% ------------------------------------------------------------------------------
\subsubsection{Integration into Development Workflow}
Import cycle detection runs as part of our continuous integration pipeline,
preventing pull requests with circular dependencies from being merged. Developers
can also run cycle detection locally before committing, receiving
immediate feedback on architectural violations. The visualization tools aid
in understanding complex dependency relationships and planning refactoring
efforts to break cycles. By enforcing acyclic dependencies automatically,
we maintain architectural integrity as the codebase evolves.

% ==============================================================================
\subsection{Package Structure and Import Conventions}

As Python codebases scale, the organization of modules and their import
relationships becomes critical to maintainability. Without clear conventions,
imports become fragile during refactoring, circular dependencies accumulate,
and understanding module relationships becomes increasingly difficult.
We developed a package organization system that prevents these problems
through clear conventions, hierarchical structure, and automated enforcement.

% ------------------------------------------------------------------------------
\subsubsection{Goals of Packages}

The goal of creating packages is to simplify imports from clients,
hide which file contains the actual code to enable reorganization without
changing client code, organize code in related units, and avoid import
loops by enforcing that there are no import loops in any module or
among modules. For example, importing from a package:

\begin{verbatim}
import dataflow.core as dtfcore
dtfcore.ArmaGenerator(...)
\end{verbatim}

is cleaner than importing the specific file:

\begin{verbatim}
import dataflow.system.source_nodes as dtfsysonod
dtfsysonod.ArmaGenerator(...)
\end{verbatim}

% ------------------------------------------------------------------------------
\subsubsection{Rules for Imports}

We follow rules to avoid import loops:

\begin{itemize}
  \item \textbf{Code inside a package imports files directly}: Within
    \texttt{datapull/common/data/client/}, code imports sibling files
    as \texttt{import datapull.common.data.client.abstract\_im\_clients
    as dcabimcl}, not the package itself. This prevents self-referential
    cycles.

  \item \textbf{Code from other packages imports the package}: External
    code imports \texttt{import datapull.common.data.client as dcdc}
    and uses \texttt{dcdc.AbstractImClient(...)}, not internal files.
    This hides implementation details and enables internal reorganization.

  \item \textbf{No statically detectable import loops}: We disallow
    any import loop that can be detected by inspecting code without
    executing it. This guarantees no dynamic import loops exist. The
    invoke command \texttt{lint\_detect\_cycles} checks for violations
    and runs automatically in continuous integration.

  \item \textbf{Only module-level imports}: Imports must occur at module
    level, not inside functions. We do not accept using local imports
    to break import loops, as they hide dependencies from static analysis.

  \item \textbf{Controlled package creation}: We avoid creating too
    many packages since organizing and maintaining them has overhead.
    We create packages only when the benefits justify the costs.

  \item \textbf{Canonical import aliases}: Each package specifies a
    short import alias in its \texttt{\_\_init\_\_.py} file, constructed
    from package path initials and kept to eight characters or fewer.
    For example, \texttt{datapull.common.data.client} becomes
    \texttt{dcdc}, \texttt{dataflow.core} becomes \texttt{dtfcore},
    and \texttt{market\_data} becomes \texttt{mdata}. The alias is
    documented in a docstring and used consistently. Linters verify
    compliance.
\end{itemize}

\paragraph{Unit Test Imports}
Unit tests can use either package imports when testing exported functionality
or direct file imports when testing internal implementation. Both are
acceptable since tests are both external clients and validators of
implementation details. However, test files are always leaf nodes and
never import from each other. Common test code belongs in library
files like \texttt{foobar\_example.py} or \texttt{test/foobar\_test\_case.py}.

% ------------------------------------------------------------------------------
\subsubsection{Package Hierarchy and Cycle Prevention}

To prevent import cycles, we enforce that certain packages do not depend
on other packages. This creates a directed acyclic graph at the package
level.

The hierarchy consists of distinct layers: \texttt{helpers} contains
low-level utilities depending only on external libraries, \texttt{core}
provides project-specific utilities depending only on \texttt{helpers},
\texttt{dataflow} implements the computational framework depending on
\texttt{core} and \texttt{helpers}, and application packages
(\texttt{oms}, \texttt{market\_data}, \texttt{im}) implement business
logic depending only on lower layers.

Within \texttt{helpers} itself, we maintain a micro-hierarchy among
utilities to prevent cycles in foundational code. The hierarchy progresses
from \texttt{hwarnings} and \texttt{hlogging} through \texttt{hdbg},
then \texttt{hintrospection} and \texttt{hprint}, to \texttt{henv},
\texttt{hsystem}, and \texttt{hio}, finally to \texttt{hgit} which
depends on system operations. A library can import only from preceding
or same-level libraries without creating cycles.

Attempts to introduce upward dependencies (such as \texttt{helpers}
importing \texttt{core}) are rejected during code review. This disciplined
approach maintains architectural clarity as the codebase evolves.

% ------------------------------------------------------------------------------
\subsubsection{Anatomy of a Package}

Every package directory contains \texttt{\_\_init\_\_.py} serving
multiple purposes: marking the directory as a Python package, documenting
the canonical import alias in a docstring, and optionally exporting
the public interface. A typical structure:

\begin{verbatim}
"""
Market data access and processing.

Import as:

import market_data as mdata
"""

from market_data.abstract_market_data import AbstractMarketData
__all__ = ["AbstractMarketData"]
\end{verbatim}

While \texttt{import *} is generally prohibited, it is acceptable within
\texttt{\_\_init\_\_.py} files for controlling the package's public
interface.

% ------------------------------------------------------------------------------
\subsubsection{Special Cases and Enforcement}

\paragraph{Unit Test Imports}
Unit tests can use either package imports when testing public interfaces
or direct file imports when testing internal implementation. Both
perspectives have merit since tests are both external clients and
validators of implementation details. However, test files must be leaf
nodes, they never import from each other. Shared test functionality
belongs in dedicated library files like \texttt{foobar\_example.py}
or \texttt{test/foobar\_test\_case.py}.

\paragraph{Enforcement Mechanisms}
Import conventions are enforced through multiple complementary mechanisms.
Static cycle detection tools analyze the dependency graph and flag
violations. Package structure verification ensures all directories
with Python files include \texttt{\_\_init\_\_.py}. Canonical alias
checking via linters verifies documented aliases are used. Automated
CI checks prevent pull requests with violations from merging. Code
review ensures new packages fit the hierarchy and architectural boundaries
are respected.

% ------------------------------------------------------------------------------
\subsubsection{Benefits}

This structured approach delivers concrete benefits. Internal file
reorganization within packages does not break external code, enabling
continuous improvement without cascading breakages. The package hierarchy
makes layering explicit, preventing ad-hoc coupling between components.
The combination of hierarchical structure and import rules makes circular
dependencies structurally difficult. Canonical aliases and consistent
patterns reduce cognitive load, letting developers focus on writing
code rather than debating import style. As the codebase grows, the
package structure provides natural boundaries for understanding subsystems
in isolation.

The primary trade-off is upfront structure. Creating new packages
requires thoughtful consideration of hierarchy placement, interface
design, and alias documentation. However, this investment pays dividends
as maintaining discipline from the start is far easier than untangling
dependencies later.

% ==============================================================================
\subsection{Pre-Commit Hook System}
Automated checks at commit time form the first line of defense against
quality regressions. While CI/CD pipelines catch issues before merging,
pre-commit hooks provide immediate feedback before committing, reducing
context switching and accelerating iteration cycles. Our pre-commit hook
system enforces essential invariants before allowing commits to succeed.

% ------------------------------------------------------------------------------
\subsubsection{Enforced Checks}
The pre-commit hook (\texttt{pre-commit.py}) executes a series of checks
automatically when developers attempt to commit changes:
\begin{itemize}
  \item \textbf{Branch verification}: Prevents accidental commits directly
    to protected branches (e.g., master)

  \item \textbf{Author validation}: Ensures Git user name and email are properly
    configured according to organizational standards

  \item \textbf{File size limits}: Blocks commits containing files exceeding
    size thresholds to prevent repository bloat

  \item \textbf{Python compilation}: Validates that all modified Python files
    compile successfully, catching syntax errors immediately

  \item \textbf{Secret detection}: Scans for accidentally committed secrets,
    credentials, or sensitive data using gitleaks integration

  \item \textbf{Forbidden patterns}: Detects merge conflict markers and other
    problematic patterns in committed files
\end{itemize}
If any check fails, the commit is aborted and developers receive a clear
diagnostic message indicating which check failed and why. Only when all checks
pass does the commit proceed.

% ------------------------------------------------------------------------------
\subsubsection{Commit Message Enhancement}
In addition to pre-commit validation, a commit-msg hook (\texttt{commit-msg.py})
enforces commit message conventions and automatically appends metadata
about executed checks. Each commit message includes a footer documenting
which pre-commit checks were run and passed, providing an audit trail of
quality enforcement. This transparency helps during code review and debugging
by clearly indicating what validations occurred.

% ------------------------------------------------------------------------------
\subsubsection{Installation and Configuration}
Git hooks are installed automatically when developers activate the thin
environment via \texttt{setenv.sh}, ensuring consistent enforcement
across all team members. The hooks can be manually installed or removed using
\texttt{install\_hooks.py} for explicit control. While hooks can be bypassed
using \texttt{git commit --no-verify} for exceptional circumstances, this
practice is discouraged and requires conscious decision-making.
Repository administrators can disable hooks for specific projects via
\texttt{repo\_config.yaml} when necessary, though the default is
universal enablement.

% ==============================================================================
\subsection{Code Coverage Tracking and Enforcement}
Maintaining comprehensive test coverage across a growing codebase requires
more than just writing tests, it demands visibility, automation, and
enforcement. Our integration with Codecov provides a system-wide view of
test coverage, structured into fast, slow, and superslow test suites. This
setup ensures that all code paths are exercised and that test coverage regressions
are identified early and reliably. The coverage system works in concert
with the linting and pre-commit systems to form a comprehensive quality
enforcement framework.

% ==============================================================================
\subsection{Structured Coverage by Test Category}

We categorize coverage tests into three suites based on runtime and
scope:

\begin{itemize}
  \item Fast tests run frequently (e.g., daily) and provide immediate feedback
    on high-priority code paths

  \item Slow tests cover broader logic and data scenarios

  \item Superslow tests are comprehensive, long-running regressions
    executed on a weekly cadence or on-demand
\end{itemize}

Each suite produces its own coverage report, which is flagged and
uploaded independently to Codecov, enabling targeted inspection and carryforward
of data when some suites are skipped.

% ==============================================================================
\subsection{CI Integration and Workflow Behavior}

Coverage reports are generated and uploaded automatically as part of our
CI pipelines. The workflow:

\begin{itemize}
  \item Fails immediately on critical setup errors (e.g., dependency or
    configuration issues)

  \item Continues gracefully if fast or slow tests fail mid-pipeline,
    but surfaces those failures in a final gating step

  \item Treats superslow failures as critical, immediately halting the
    workflow
\end{itemize}

This behavior ensures resilience while preventing silent test
degradation.

% ==============================================================================
\subsection{Enforced Thresholds and Quality Gates}

Coverage checks are enforced at both project and patch levels:

\begin{itemize}
  \item Project-level threshold: Pull requests fail if overall coverage
    drops beyond a configured margin (e.g., \textgreater1\%)

  \item Patch-level checks: Changes are required to maintain or improve
    coverage on modified lines

  \item Flags and branches: Checks are scoped per test suite and only
    enforced on critical branches
\end{itemize}

Together, these gates maintain coverage integrity while avoiding noise from
unrelated code paths.

% ==============================================================================
\subsection{Visibility and Developer Experience}

Codecov is integrated tightly into the developer workflow:

\begin{itemize}
  \item PRs show inline coverage status and file-level diffs

  \item Optional summary comments detail total coverage, changes, and affected
    files

  \item Reports can be viewed in Codecov's UI or served locally as HTML

  \item Carryforward settings retain historical data when full test suites
    aren't executed
\end{itemize}

Developers can also generate and inspect local reports for any test suite
using standard coverage commands.

% ==============================================================================
\subsection{Best Practices and Operational Consistency}

To ensure effective usage:

\begin{itemize}
  \item Coverage is always uploaded---even if tests fail---ensuring no blind
    spots

  \item Developers are encouraged to monitor coverage deltas in PRs

  \item The system defaults to global configuration, but supports fine-tuning
    via repo-specific overrides

  \item Weekly reviews of coverage trends and flags help spot regressions
    and low-tested areas
\end{itemize}

% ##############################################################################
\section{Unit test framework}

% Related documentation files:
% - helpers_root/docs/tools/unit_test/all.write_unit_tests.how_to_guide.md
%   Core guide to writing unit tests with hunitest framework
% - helpers_root/docs/tools/unit_test/all.run_unit_tests.how_to_guide.md
%   Running and organizing unit tests
% - helpers_root/docs/tools/all.code_coverage.how_to_guide.md
%   Measuring test coverage
% - helpers_root/docs/tools/all.codecov.how_to_guide.md
%   Coverage reporting and enforcement
% - helpers_root/docs/tools/notebooks/all.run_jupyter_notebook.how_to_guide.md
%   Testing Jupyter notebooks
% - helpers_root/docs/tools/notebooks/all.jupyter_notebook.how_to_guide.md
%   Working with notebooks in tests
% - helpers_root/docs/build_system/all.pytest_allure.explanation.md
%   Test reporting with Allure
% - helpers_root/docs/build_system/all.pytest_allure.how_to_guide.md
%   Setting up Allure for test visualization
% - docs/oms/broker/ck.generate_broker_test_data.how_to_guide.md
%   Example of generating test data for complex systems
% - docs/oms/broker/ck.replayed_ccxt_exchange.explanation.md
%   Mocking external dependencies for testing

% ==============================================================================
\subsection{Summary}
% TODO(gp_ai3): The text should correspond to the topics here and convert to latex

%Our unit testing framework combines principled philosophy with practical
%tooling.
%
%- By inheriting from \texttt{hunitest.TestCase}, tests gain
%access to enhanced assertions, golden file testing, and directory
%management utilities.
%
%- Categorizing tests by speed enables fast feedback cycles
%without sacrificing comprehensive regression testing.
%
%- The opinionated
%approach to mocking—mock only external dependencies—keeps tests focused
%on behavior and maintainable over time. 
%
%- Coverage measurement and enforcement
%ensure quality gates without creating busywork.
%
%- check_string vs assert_equal
%  - Support for fuzzy testing
%
%- Test organization and conventions
%
%- Mocking conventions
%
%- Special support for
%Jupyter notebooks extends testing to interactive, exploratory workflows.
%This testing infrastructure is foundational to our development system. It
%enables confident refactoring, documents expected behavior, and catches regressions
%early.
%
%- Test coverage
%
%- The consistency of patterns and conventions reduces cognitive
%load, allowing developers to focus on writing tests that matter rather than
%debating how to structure them.

% Different levels of testing
% Where to save tests

% ==============================================================================
\subsection{Testing Philosophy and Motivation}
Testing needs to be a first-class concern in software development.

We view unit tests as
executable specifications of behavior. Each test documents what the
system should do, provides a safety net for refactoring, and serves as
living documentation for developers. Good tests improve code design by forcing
clear interfaces and reducing coupling between components.

% ==============================================================================
\subsection{The hunitest.TestCase Framework}
At the core of our testing infrastructure is \texttt{hunitest.TestCase},
a custom test base class that extends Python's standard \texttt{unittest}
framework with specialized functionality for our development workflows.
Every test class in our codebase inherits from \texttt{hunitest.TestCase},
ensuring consistent testing patterns and access to enhanced testing
utilities. The framework provides several categories of functionality:
\begin{itemize}
  \item \textbf{Enhanced assertions}: Methods like \texttt{assert\_equal()}
    that provide superior diff output compared to standard \texttt{assertEqual()},
    especially for large strings and data structures

  \item \textbf{Golden file testing}: The \texttt{check\_string()}
    method for regression testing against frozen reference outputs
    stored on disk

  \item \textbf{Directory management}: Automatic creation and cleanup of
    test directories for inputs, outputs, and scratch space

  \item \textbf{Text processing}: Utilities for fuzzy matching,
    whitespace normalization, and memory address purification in test assertions

  \item \textbf{Notebook execution}: Specialized support for testing Jupyter
    notebooks as executable artifacts
\end{itemize}

% ------------------------------------------------------------------------------
\subsubsection{Basic Test Structure}
Every test method follows a consistent three-section structure that
makes tests self-documenting and easy to understand:

\begin{lstlisting}[language=Python]
def test_something(self) -> None:
  """
  Brief description of what this tests.
  """
  # Prepare inputs.
  input_data = create_test_data()
  # Run test.
  actual = function_under_test(input_data)
  # Check outputs.
  expected = "expected_result"
  self.assert_equal(str(actual), str(expected))
\end{lstlisting}

This structure enforces clarity: setup is separate from execution, which
is separate from verification. When a test fails, developers can immediately
identify which phase failed and why. The mandatory docstring and type
hints further improve readability and maintainability.

% ==============================================================================
\subsection{Golden File Testing with check\_string}
For complex outputs or large data structures, comparing against inline expected
values becomes impractical. Golden file testing solves this by storing
reference outputs in versioned files alongside the test code.
The \texttt{check\_string()} method implements this pattern.

% ------------------------------------------------------------------------------
\subsubsection{How check\_string Works}
When a test calls \texttt{self.check\_string(actual)}, the framework:
\begin{enumerate}
  \item Converts the actual output to a string representation

  \item Computes a path based on the test class and method name (e.g.,
    \texttt{test/outcomes/TestFoo1.test\_bar1/output/test.txt})

  \item Compares the actual output against the stored golden file

  \item Reports any differences using sophisticated diff algorithms

  \item Provides a command to update the golden file if the change is intentional
\end{enumerate}

This approach has several advantages over inline assertions:
\begin{itemize}
  \item Large outputs don't bloat the test code

  \item Diffs are easier to review in version control

  \item Tests remain readable even for complex outputs

  \item Updating expectations is a single command: \texttt{pytest --update\_outcomes}

  \item Historical reference outputs are preserved in git history
\end{itemize}

% ------------------------------------------------------------------------------
\subsubsection{When to Use check\_string vs assert\_equal}
The framework provides clear guidance:
\begin{itemize}
  \item Use \texttt{assert\_equal()} for: Simple comparisons (e.g.,
    \texttt{1 + 1 == 2}), small outputs (< 50 lines), values that should
    never change

  \item Use \texttt{check\_string()} for: Large outputs (> 50 lines),
    complex data structures, outputs that may evolve with the system,
    regression testing of behavior
\end{itemize}
Both methods support optional parameters for flexible comparison: \texttt{fuzzy\_match=True}
for whitespace-insensitive comparison, \texttt{purify\_text=True} for
stripping memory addresses and timestamps, and \texttt{dedent=True} for
automatic indentation normalization.

% ==============================================================================
\subsection{Test Organization and Conventions}
Consistent test organization reduces cognitive overhead and makes the codebase
easier to navigate. Our conventions cover naming, file placement, and
structural patterns.

% ------------------------------------------------------------------------------
\subsubsection{File and Directory Structure}
Tests live in \texttt{test/} directories adjacent to the code they test.
For a module \texttt{foo/bar/module.py}, the test file is \texttt{foo/bar/test/test\_module.py}.
Test artifacts (golden files, input data) are stored in subdirectories
following the pattern \texttt{test/TestClassName.test\_method\_name/}.
This co-location ensures tests are easy to find and maintain. When code moves,
tests move with it. When code is deleted, tests are deleted too.

% ------------------------------------------------------------------------------
\subsubsection{Naming Conventions}
Test class names clearly indicate what they test:
\begin{itemize}
  \item For a class \texttt{FooBar}: \texttt{TestFooBar1}

  \item For a function \texttt{process\_data()}: \texttt{Test\_process\_data1}

  \item For a method \texttt{FooBar.calculate()}: test method \texttt{TestFooBar1.test\_calculate1()}
\end{itemize}
Numbers (e.g., \texttt{TestFooBar1}, \texttt{test\_calculate1()}) are used
even for the first test to facilitate adding additional test classes or methods
later without renaming files and directories. Test methods use
descriptive names when there are few variations (\texttt{test\_calculate\_with\_empty\_input()})
or numbered names when testing multiple similar scenarios (\texttt{test\_calculate1()},
\texttt{test\_calculate2()}, \texttt{test\_calculate3()}).

% ------------------------------------------------------------------------------
\subsubsection{Helper Methods and DRY Principle}
When multiple test methods follow the same pattern with different inputs,
helper methods eliminate repetition:

\begin{lstlisting}[language=Python]
class TestFooBar1(hunitest.TestCase):
  def _helper(self, input_val, expected):
      # Run test.
      actual = foo_bar(input_val)
      # Check outputs.
      self.assert_equal(str(actual), str(expected))

  def test1(self) -> None:
      self._helper(input_val=1, expected="result1")

  def test2(self) -> None:
      self._helper(input_val=2, expected="result2")
\end{lstlisting}

This keeps tests DRY (Don't Repeat Yourself) while maintaining clarity about
what distinguishes each test case.

% ==============================================================================
\subsection{Test Categorization and Execution}
Tests are categorized by execution time to enable fast feedback cycles. The
framework defines three categories using pytest markers:

\begin{table}[h]
\centering
\begin{tabular}{|l|p{5cm}|p{5cm}|}
\hline
\textbf{Test Category} & \textbf{Examples of Tests} & \textbf{When to Run} \\
\hline
\textbf{Fast tests} (< 5 sec per test class) & Core functionality, business logic, unit-level tests & Run before every commit and in CI on every pull request \\
\hline
\textbf{Slow tests} (< 20 sec) & More complex scenarios, integration-level tests, tests involving external APIs (mocked) & Run daily and before releases \\
\hline
\textbf{Superslow tests} (< 5 min) & End-to-end workflows, production simulations, comprehensive regression tests & Run weekly or on-demand \\
\hline
\end{tabular}
\caption{Test categorization by execution time}
\label{tab:test_categories}
\end{table}

% ------------------------------------------------------------------------------
\subsubsection{Running Tests with invoke}
The invoke task system provides convenient commands for running each
category:

\begin{lstlisting}[language=bash]
# Run Fast Tests Only
invoke run_fast_tests
# Run Slow Tests
invoke run_slow_tests
# Run Superslow Tests
invoke run_superslow_tests
# Run with Coverage
invoke run_fast_tests --coverage
# Run Specific Test
invoke run_fast_tests --pytest-opts="path/to/test.py::TestClass::test_method"
\end{lstlisting}
Tests always execute inside Docker containers to ensure environmental consistency.
The framework automatically handles pytest options, timeout configuration,
and result reporting.

% ------------------------------------------------------------------------------
\subsubsection{Timeout and Retry Behavior}
Tests are subject to timeouts based on their category to prevent runaway
tests from blocking CI pipelines. The \texttt{pytest-timeout} plugin enforces
these limits. Additionally, \texttt{pytest-rerunfailures} automatically
retries flaky tests: fast tests are retried twice, slow and superslow
tests once. This balances reliability with execution time.

% ==============================================================================
\subsection{Mocking Philosophy and Practice}
Our approach to mocking is pragmatic and opinionated: mock only external
dependencies, never mock internal code. This philosophy stems from our end-to-end
testing preference and our focus on behavioral correctness.

% ------------------------------------------------------------------------------
\subsubsection{Mock Only External Dependencies}
We mock interactions with systems outside our control:
\begin{itemize}
  \item Third-party APIs (e.g., CCXT cryptocurrency exchange library)

  \item Cloud infrastructure (S3, AWS Secrets Manager, databases)

  \item External services (GitHub API, OpenAI, data providers)
\end{itemize}
For example, when testing code that fetches data from a cryptocurrency
exchange via CCXT, we mock the CCXT library, not our wrapper around it:

\begin{lstlisting}[language=Python]
@umock.patch.object(module.ccxt, "binance")
def test_fetch_ohlcv(self, mock_binance):
  mock_exchange = umock.Mock()
  mock_binance.return_value = mock_exchange
  mock_exchange.fetch_ohlcv.return_value = test_data
  # Now test our code's behavior
\end{lstlisting}
This ensures we're actually testing our code's logic, not the mock
itself.

% ------------------------------------------------------------------------------
\subsubsection{Do Not Mock Internal Code}
We never mock code within our own repositories. Doing so creates
maintenance burdens and tests implementation details rather than
behavior. If internal code is difficult to test without mocking, that's
a signal to refactor for better testability. This approach means tests exercise
real code paths, catching integration issues that mocking would hide. It
also makes refactoring safer: as long as the public interface remains
stable, internal restructuring doesn't require updating tests.

% ------------------------------------------------------------------------------
\subsubsection{Shared Mock Setup}
For test classes that require common mocks, we use \texttt{set\_up\_test()}
and \texttt{tear\_down\_test()} with pytest fixtures:

\begin{lstlisting}[language=Python]
class TestExchangeClient1(hunitest.TestCase):
  # Create patches as class variables
  ccxt_patch = umock.patch.object(module, "ccxt")
  secret_patch = umock.patch.object(module.hsecret, "get_secret")
  @pytest.fixture(autouse=True)
  def setup_teardown_test(self):
      self.set_up_test()
      yield
      self.tear_down_test()
  def set_up_test(self) -> None:
      self.ccxt_mock = self.ccxt_patch.start()
      self.secret_mock = self.secret_patch.start()
      self.secret_mock.return_value = {"key": "test"}
  def tear_down_test(self) -> None:
      self.ccxt_patch.stop()
      self.secret_patch.stop()
\end{lstlisting}

This pattern ensures mocks are consistently configured across all test
methods and properly cleaned up after each test.

% ==============================================================================
\subsection{Test Coverage Measurement and Enforcement}
Code coverage is measured using the \texttt{coverage.py} library
integrated with pytest via \texttt{pytest-cov}. Coverage is tracked
separately for fast, slow, and superslow test suites, providing visibility
into which code paths are exercised by which test categories.

% ------------------------------------------------------------------------------
\subsubsection{Running Coverage Analysis}
Coverage analysis is initiated with the \texttt{--coverage} flag:

\begin{lstlisting}[language=bash]
# Generate Coverage for Fast Tests
invoke run_fast_tests --coverage
# Generate Coverage for Specific Module
invoke run_fast_tests --coverage --pytest-opts="oms/"
# Combine Coverage From Multiple Runs
docker> coverage combine .coverage_fast .coverage_slow
docker> coverage report --include="oms/*" --omit="*/test_*.py"
\end{lstlisting}

The framework generates both terminal reports and browsable HTML
coverage reports (\texttt{htmlcov/index.html}) showing line-by-line coverage
with color-coded indicators for covered, uncovered, and partially covered
lines.

% ------------------------------------------------------------------------------
\subsubsection{Coverage Integration with CI/CD}
Coverage reports are uploaded to Codecov as part of our CI/CD pipelines.
Each test suite (fast, slow, superslow) produces a separate coverage flag,
enabling targeted analysis. Pull requests are gated on coverage thresholds:
project-level coverage must not decrease by more than 1\%, and modified
lines must maintain or improve coverage. This prevents coverage regression
while allowing flexibility for exploratory development.

% ------------------------------------------------------------------------------
\subsubsection{Interpreting Coverage Metrics}
We focus on behavioral coverage, not line coverage. The goal is not 100\%
line coverage but comprehensive testing of observable behaviors and critical
code paths. Some lines, like defensive assertions and error handling for
impossible states, may remain uncovered by design. Coverage reports guide,
rather than dictate, testing efforts. Low coverage on a module signals
potential risk, but the appropriate response depends on context:
critical trading logic demands exhaustive testing, while simple utility functions
may not.

% ==============================================================================
\subsection{Testing Jupyter Notebooks}
Jupyter notebooks present unique testing challenges. They combine code,
narrative, and visualizations in a format designed for interactive exploration,
not automated testing. Our framework treats notebooks as executable artifacts
that must be validated for correctness and reproducibility.

% ------------------------------------------------------------------------------
\subsubsection{Notebook Execution Framework}
Notebooks are tested by executing them in a controlled environment using
\texttt{run\_notebook.py}, a wrapper around \texttt{papermill} that injects
parameters, captures output, and validates results:

\begin{lstlisting}[language=bash]
# Parameterized Notebook Execution
/app/dev_scripts/notebooks/run_notebook.py \
--notebook path/to/notebook.ipynb \
--config_builder "config_builder_function" \
--dst_dir /tmp/results
\end{lstlisting}

This approach ensures notebooks execute without errors, produce expected
outputs, and remain compatible with current code versions.

% ------------------------------------------------------------------------------
\subsubsection{Notebook Test Organization}
Notebook tests follow the same naming conventions as regular tests. For
a notebook \texttt{analysis.ipynb}, the test is \texttt{test\_analysis\_notebook.py}.
The test class typically has a single test method that executes the
notebook and validates key outputs or side effects. Tests use golden file
testing to detect unintended changes in notebook output. When notebooks
generate plots, tables, or metrics, these are captured and compared against
reference versions stored in test outcomes.

% ------------------------------------------------------------------------------
\subsubsection{Debugging Failing Notebook Tests}
When a notebook test fails, the debugging workflow is:
\begin{enumerate}
  \item Run the failing test with \texttt{-s --dbg} to get detailed logs

  \item Extract the \texttt{run\_notebook.py} command from the logs

  \item Re-run with \texttt{--no\_suppress\_output} to see full notebook
    output

  \item Identify the failing cell and root cause from the detailed error
    report
\end{enumerate}

This workflow provides full visibility into notebook execution while
keeping test output concise during normal runs.

% ##############################################################################
\section{invoke workflows}

% Related documentation files:
% - helpers_root/docs/tools/all.invoke_workflows.how_to_guide.md
%   Core guide to using invoke for task automation
% - helpers_root/docs/tools/all.invoke_git_branch_copy.how_to_guide.md
%   Specific invoke task example for git operations
% - helpers_root/docs/tools/all.invoke_bash_print_tree.how_to_guide.md
%   Invoke task for directory visualization
% - docs/infra/ck.kaizen_infra_script.how_to_guide.md
%   Infrastructure invoke tasks
% - docs/infra/ck.kaizen_infra_python_script.how_to_guide.md
%   Python-based invoke workflows
% - docs/infra/ck.kaizen_infra_python_script.reference.md
%   Reference for infra invoke tasks
% - helpers_root/docs/onboarding/all.development.how_to_guide.md
%   Common invoke commands for daily development

% Invoke Task Automation: Centralized task registry using Python Invoke for common development operations (testing, linting, deployment).

% TODO(ai_gp2): Add Motivation

% TODO(ai_gp2): Add Summary

% ##############################################################################
\section{Code Repo Automation}

% Related documentation files:
% - helpers_root/docs/tools/dev_system/all.synchronize_gh_issue_labels.how_to_guide.md
%   Automated label synchronization
% - helpers_root/docs/tools/dev_system/all.synchronize_gh_settings.how_to_guide.md
%   Repository settings synchronization
% - helpers_root/docs/tools/all.sync_gh_projects.how_to_guide.md
%   GitHub Projects synchronization
% - helpers_root/docs/tools/git/all.invite_gh_contributors.explanation.md
%   Automating contributor management
% - helpers_root/docs/tools/git/all.invite_gh_contributors.how_to_guide.md
%   How to manage contributors programmatically
% - helpers_root/docs/work_organization/all.github_projects_process.reference.md
%   GitHub Projects structure and automation
% - helpers_root/docs/work_organization/all.pr_workflow.explanation.md
%   Pull request automation workflow
% - helpers_root/docs/work_organization/all.issue_workflow.explanation.md
%   Issue lifecycle automation

% ==============================================================================
\subsection{Motivation}

As organizations grow, so does the surface area of their code repo
ecosystem---more repositories, more contributors, and more operational
complexity. Without strong conventions and enforcement, labels diverge
across projects, and team boards lose structural consistency. This leads
to disjointed workflows, unclear ownership, and fragmented planning. To mitigate
this, we introduced a declarative automation system for GitHub metadata.
It enables centralized definitions for labels and project structures and
propagates them across the organization in a reproducible and scalable way.

We settle on git as version control system and GitHub as ...

Multiple alternative solutions are available that would 

% ==============================================================================
\subsection{Summary of Solutions}

- Make it easy to create composable Git modules for reuse and access control
- Ensure that each Git repo follows the same standards (e.g., unified labels,
  configuring GitHub projects)
- A declarative approach to repository settings

% ==============================================================================
\subsection{Standardized Label Infrastructure}

Issue labels form the foundation for triage, workflows, and reporting. At
scale, inconsistent naming, coloring, or descriptions can fragment
automation and reduce clarity. To address this, we maintain a
centralized manifest that defines the entire organization's label
taxonomy. Each label includes a name, color code, and description.
Repositories synchronize against this manifest using a containerized
process that:

\begin{itemize}
  \item Creates missing labels using manifest definitions

  \item Updates outdated labels (e.g., if the description or color
    changes)

  \item Backs up existing labels before applying changes

  \item Supports dry-run execution for visibility and confidence

  \item Optionally prunes unused labels not defined in the manifest
\end{itemize}

This guarantees all repositories speak a consistent language for issues and
pull requests, improving developer experience and enabling reliable
automation.

% ==============================================================================
\subsection{Template-Based Project Synchronization}

GitHub Projects (Beta) are a powerful tool for planning and tracking,
but manually configuring project fields and views across teams is
tedious and error-prone. To streamline setup and reduce drift, we
introduced a project templating system. Project metadata, such as fields
and views is defined in a canonical source project and synced into
destination projects. The system:

\begin{itemize}
  \item Adds missing global fields to ensure functional parity across projects

  \item Logs discrepancies in view structure (e.g., missing ``Backlog''
    or ``Current sprint'' views)

  \item Preserves existing configurations to avoid accidental data loss

  \item Operates in a dry-run mode to preview proposed changes safely
\end{itemize}

Due to current API limitations, view creation, layout configuration, and
field ordering cannot yet be automated. However, warnings are logged for
manual follow-up, and the system is designed to evolve as GitHub expands
its GraphQL support.

% ==============================================================================
\subsection{Design Principles and Workflow Behavior}

% TODO(ai_gp2): Merge these two things

These automation tools follow several core principles:

\begin{itemize}
  \item Declarative Configuration: Metadata is defined in structured
    YAML files stored in version control, serving as the source of truth

  \item Reproducible Execution: All sync operations run in isolated
    Docker containers, ensuring consistent behavior across machines and CI
    pipelines

  \item Non-Destructive by Default: No changes are applied unless
    explicitly confirmed, and current states can be backed up for rollback

  \item Extensibility: The architecture is modular and designed to
    incorporate future GitHub API enhancements, including view creation and
    layout syncing
\end{itemize}

The workflows support both interactive execution and automated pipelines,
allowing integration into CI/CD systems or local tooling as needed.

% ==============================================================================
\subsection{Declarative Repository Settings Synchronization}

As code base footprint grows, so has the need to manage repository
configuration such as metadata fields, merge policies, and branch protection
rules in a reliable and consistent manner across dozens of projects. Manual
configuration inevitably leads to drift, accidental misconfiguration, and
wasted effort during onboarding or project setup. To address this, we
introduced a \textit{declarative} synchronization system for repository
settings, inspired by the principles of infrastructure-as-code. Repository
policies—including key settings (description, default branch, visibility, merge
strategies, etc.) and fine-grained branch protection rules are defined
centrally in a YAML manifest. This manifest serves as the single source of
truth for desired state. A containerized executable can \textbf{export} the
current settings of any repository for inspection, backup, or review, and can
\textbf{apply} (sync) the manifest to one or more repositories, ensuring
alignment with organizational standards. The system also supports a dry-run
mode for safe preview of all proposed changes.

This approach delivers several benefits:
\begin{itemize}
  \item \textbf{Consistency:} All repositories adhere to the same
    baseline policies, reducing configuration drift and error.

  \item \textbf{Auditability:} Changes to settings are reviewed and
    version-controlled like code, supporting transparent change
    management.

  \item \textbf{Scalability:} New projects can be onboarded with a
    single command, inheriting standardized policies from day one.

  \item \textbf{Recoverability:} Before each sync, the system backs up the
    existing configuration, enabling rollbacks if needed.

  \item \textbf{Confidence:} Dry-run and logging capabilities provide
    full visibility into what changes would be made, supporting both
    local experimentation and safe automation in CI/CD pipelines.
\end{itemize}

% ##############################################################################
\section{CI/CD Automation}

% Related documentation files:
% - helpers_root/docs/work_organization/all.buildmeister.how_to_guide.md
%   Core Buildmeister role and responsibilities
% - helpers_root/docs/work_organization/all.inframeister.how_to_guide.md
%   Infrastructure-focused CI/CD monitoring
% - helpers_root/docs/work_organization/all.datapull_DAGmeister.how_to_guide.md
%   Data pipeline CI/CD monitoring
% - helpers_root/docs/build_system/all.pytest_allure.explanation.md
%   Allure test reporting in CI/CD
% - helpers_root/docs/build_system/all.pytest_allure.how_to_guide.md
%   Setting up Allure reports
% - helpers_root/docs/build_system/all.gitleaks_workflow.explanation.md
%   Gitleaks secret detection in CI/CD
% - helpers_root/docs/build_system/all.linter_gh_workflow.explanation.md
%   Linter automation in GitHub Actions
% - helpers_root/docs/build_system/all.semgrep_workflow.explanation.md
%   Semgrep security scanning workflow
% - docs/build/ck.github_actions.explanation.md
%   GitHub Actions CI/CD setup
% - docs/build/ck.gh_workflows.explanation.md
%   Workflow definitions and structure
% - docs/build/ck.gh_workflows.reference.md
%   Complete workflow reference
% - docs/build/ck.setup_github_actions.how_to_guide.md
%   Initial GitHub Actions setup
% - docs/build/ck.debug_github_actions.how_to_guide.md
%   Debugging CI/CD failures
% - docs/build/ck.github_runners.explanation.md
%   GitHub runners configuration
% - docs/build/ck.setup_slack_notification_for_actions.how_to_guide.md
%   Build notification setup

% ==============================================================================
\subsection{Motivation}

Automated test pipelines are essential, but without accountability, they
often fall into disrepair. The Buildmeister routine introduces a
rotating, human-in-the-loop system designed to enforce green builds,
identify root causes, and ensure high-quality CI/CD hygiene. This
mechanism aligns technical execution with team responsibility, fostering
a culture of operational ownership and co-operation.

% ==============================================================================
\subsection{Why It Matters}

The Buildmeister is a system of shared accountability. It transforms test
stability from an abstract ideal into a daily operational habit, backed by
clear expectations, defined processes, and human enforcement. By combining
automation with ownership, the team achieves sustainable reliability in a
complex, multi-repo ecosystem.

% ==============================================================================
\subsection{Core Responsibilities}

The Buildmeister is a rotating role assigned to a team member every one or two weeks. Their primary duties are:

\begin{itemize}
  \item Monitor build health daily via the Buildmeister Dashboard

  \item Investigate failures and ensure corresponding GitHub Issues are filed
    promptly

  \item Push responsible team members to fix or revert breaking code

  \item Maintain test quality by analyzing trends in reports

  \item Document breakage through a structured post-mortem log
\end{itemize}

The Buildmeister ensures the policy of: ``Fix it or revert within one hour.''

% ==============================================================================
\subsection{Handover and Daily Reporting}

The routine begins each day with a status email to the team detailing:

\begin{itemize}
  \item Overall build status (green/red)

  \item Failing test names and owners

  \item GitHub issue references

  \item Expected resolution timelines

  \item A screenshot of the Buildmeister dashboard
\end{itemize}

At the end of each rotation, the outgoing Buildmeister must confirm handover
by receiving a reply from the incoming one, ensuring continuity and awareness.

% ==============================================================================
\subsection{Workflow in Practice}

When a build breaks:

\begin{itemize}
  \item The team is alerted via Slack (\#build-notifications) through
    our GitHub Actions bot

  \item The Buildmeister triages the issue:

    \begin{itemize}
      \item Quickly reruns or replicates the failed tests if uncertain

      \item Blames commits to identify the responsible party

      \item Notifies the team and files a structured GitHub Issue
    \end{itemize}

  \item All information including test names, logs, responsible engineer
    are transparently shared and tracked
\end{itemize}

If the issue is not resolved within one hour, the Buildmeister must escalate
and disable the test with explicit owner consent or revert the offending change.

% ==============================================================================
\subsection{Tools and Analysis}

% ------------------------------------------------------------------------------
\subsubsection{Buildmeister Dashboard}

A centralized UI provides a real-time view of all builds across repos and
branches. It is the Buildmeister's daily launchpad.

% ------------------------------------------------------------------------------
\subsubsection{Allure Reports}

\begin{itemize}
  \item Every week, the Buildmeister reviews trends in skipped/failing
    tests, duration anomalies, and retry spikes

  \item This process:

    \begin{itemize}
      \item Surfaces hidden test instability

      \item Provides historical context to new breaks

      \item Enables preventive action before regressions cascade
    \end{itemize}
\end{itemize}

% ------------------------------------------------------------------------------
\subsubsection{Post-Mortem Log}

Every build break is logged in a shared spreadsheet, capturing:

\begin{itemize}
  \item Repo and test type

  \item Link to the failing GitHub run

  \item Root cause

  \item Owner and fix timeline

  \item Whether the issue was fixed or test was disabled
\end{itemize}

This living record forms the basis for failure mode analysis and future automation
improvements.


  % ############################################################################
  % ############################################################################
\section{Containerized development workflows}

% Related documentation files:
% - helpers_root/docs/tools/docker/all.docker.tutorial.md
%   Comprehensive Docker tutorial covering development workflows
% - helpers_root/docs/tools/docker/all.dockerized_flow.explanation.md
%   Explains the containerized development flow
% - helpers_root/docs/tools/docker/all.dind_and_sibling_containers.how_to_guide.md
%   Details on Docker-in-Docker and sibling container patterns
% - helpers_root/docs/tools/docker/all.develop_in_helpers.how_to_guide.md
%   Developing within containerized environments
% - helpers_root/docs/tools/thin_environment/all.thin_environment.reference.md
%   Reference for the thin environment concept
% - helpers_root/docs/tools/thin_environment/all.gh_thin_env_dependencies.how_to_guide.md
%   Managing thin environment dependencies
% - helpers_root/docs/tools/thin_environment/all.gh_and_thin_env_requirements.reference.md
%   Requirements and specifications for thin environments
% - helpers_root/docs/tools/dev_system/all.dealing_with_missing_dependencies.how_to_guide.md
%   Troubleshooting container dependency issues

A ``thin environment'' contains the minimal dependency setup to bootstrap
development and deployment using ``runnable directories" 
across diverse systems (servers, laptops, CI/CD). Runnable directories are a
Docker-based system that ensures reproducible environments across development,
testing, and production stages, with support for both children and sibling
container patterns.

% TODO(gp):

% ############################################################################
\section{The \texttt{//helpers} submodule}

% Related documentation files:
% - helpers_root/docs/tools/dev_system/all.create_a_super_repo_with_helpers.how_to_guide.md
%   Explains how to create a repository that includes helpers as submodule
% - helpers_root/docs/tools/dev_system/all.integrate_repos.how_to_guide.md
%   Guide on integrating multiple repositories using helpers
% - helpers_root/docs/tools/dev_system/all.replace_common_files_with_script_links.md
%   Details the symbolic link strategy for sharing common files
% - helpers_root/docs/tools/git/all.integrate_helpers.how_to_guide.md
%   Instructions for integrating helpers submodule into repositories
% - helpers_root/docs/onboarding/all.development_setup.how_to_guide.md
%   Initial setup including helpers submodule installation

Helpers is a foundational Python utility library that serves as the core
infrastructure for an ecosystem of software development repositoriesIt's
designed to be integrated as a git submodule into ``super-repos" that extend
its functionality.

It contains all the code for all the functionalities described in this
paper.

% ############################################################################
\section{Runnable directory architecture}

% Related documentation files:
% - helpers_root/docs/tools/dev_system/all.runnable_repo.reference.md
%   Core reference for runnable directory architecture
% - helpers_root/docs/tools/dev_system/all.create_a_runnable_dir.how_to_guide.md
%   Step-by-step guide to creating a runnable directory
% - helpers_root/docs/tools/dev_system/all.devops_docker.reference.md
%   Reference for DevOps setup in runnable directories
% - helpers_root/docs/tools/dev_system/all.devops_docker.how_to_guide.md
%   How to set up DevOps workflows for runnable directories
% - docs/kaizenflow/all.install_helpers.how_to_guide.md
%   Installing and using helpers across multiple runnable directories
% - infra/docs/operations/ck.infra_runnable_dir.how_to_guide.md
%   Infrastructure setup for runnable directories

% TODO(gp): Merge in the full paper

% ============================================================================
\subsection{Problem}

Software development teams face a fundamental organizational dilemma when
structuring their codebases. The choice between monorepos and multi-repos
presents a difficult trade-off with no universally satisfactory solution.

Monorepos provide environment consistency and simplified version control,
but struggle with scalability.

Multi-repos offer modularity and independent development lifecycles, but
introduce synchronization issues between repositories.

The core challenge is achieving both modularity and consistency
simultaneously---allowing teams to work independently on separate components
while ensuring all components operate in compatible, reproducible environments.
Traditional approaches force teams to sacrifice one for the other.

% ============================================================================
\subsubsection{The monorepo approach}

% TODO(gp): Fix the references

The monorepo approach involves storing all code for multiple applications
within a single repository. This strategy has been popularized by
large tech companies like Google[2], Meta[3], Microsoft[4] and Uber[5],
proving that even codebases with billions of lines of code can be effectively
managed in a single repository. The key benefits of this approach
include:

\begin{itemize}
  \item Consistency in environment: with everything housed in one
    repository, there's no risk of projects becoming incompatible due to
    conflicting versions of third-party packages.

  \item Simplified version control: there is a single commit history, which
    makes it easy to track and, if needed, revert changes globally.

  \item Reduced coordination overhead: developers work within the same
    repository, with easy access to all code, shared knowledge, tools and
    consistent coding standards.
\end{itemize}

However, as monorepo setups scale, users often face significant challenges.
A major downside is long CI/CD build times, as even small changes can trigger
massive rebuilds and tests throughout the entire codebase. To cope with this,
extra tooling, such as [Buck](https://buck2.build/) or
[Bazel](https://bazel.build/), must be configured, adding complexity to
workflows. Even something as simple as searching and browsing the code becomes
more difficult, often requiring specialized tools and IDE plug-ins.

Additionally, when everything is located in one place, it is harder to
separate concerns and maintain clear boundaries between projects. Managing
permissions also becomes more difficult when only selected developers should
have access to specific parts of the codebase.

% ============================================================================
\subsubsection{The monorepo approach}

The multi-repo approach involves splitting code across several
repositories, with each one dedicated to a specific module or service.
This modularity allows teams to work independently on different parts
of a system, making it easier to manage changes and releases for individual
components. Each repository can evolve at its own pace, and developers
can focus on smaller, more manageable codebases.

However, the multi-repo strategy comes with its own set of challenges,
particularly when it comes to managing dependencies and ensuring
version compatibility across repositories. For instance, different
repositories might rely on two different versions of a third-party
package, or even conflicting packages, making synchronization complex
or, in some cases, nearly impossible. In general, propagating changes
from one repository to another requires careful coordination. Tools like
[Jenkins](https://www.jenkins.io/) and [GitHub Actions](https://github.com/features/actions)
help streamline CI/CD pipelines, but they often struggle when dealing with
heterogeneous environments.

% ============================================================================
\subsection{Solution: Runnable directories}

An ideal strategy would combine the best of both worlds:

\begin{itemize}
  \item The modularity of multi-repos, to keep the codebase scalable
    and simplify day-to-day development processes.

  \item The environment consistency of monorepos, to avoid
    synchronization issues and prevent errors that arise from
    executing code in misaligned environments.
\end{itemize}

Both are achieved through the hybrid approach proposed in this paper,
which will be discussed in Section 3.

\begin{itemize}
  \item This section describes the design principles in our approach to
    create Git repos that contain code that can be:

    \begin{itemize}
      \item Composed through different Git sub-module

      \item Tested, built, run, and released (on a per-directory basis
        or not)
    \end{itemize}

  \item The technologies that this approach relies on are:

    \begin{itemize}
      \item Git for source control

      \item Python virtual environment and \texttt{poetry} (or similar)
        to control Python packages

      \item \texttt{pytest} for unit and end-to-end testing

      \item Docker for managing containers
    \end{itemize}

  \item The approach described in this paper is not strictly dependent
    of the specific package (e.g., \texttt{poetry} can be replaced by \texttt{Conda}
    or another package manager)
\end{itemize}

% ============================================================================
\subsection{Design goals}

The proposed development system supports the following functionalities

% ----------------------------------------------------------------------------
\subsubsection{Development functionalities}

\begin{itemize}
  \item Support composing code using a GitHub sub-module approach

  \item Make it easy to add the development tool chain to a ``new project''
    by simply adding the Git sub-module \texttt{//helpers} to the
    project

  \item Create complex workflows (e.g., for dev and devops
    functionalities) using makefile-like tools based on Python \texttt{invoke}
    package

  \item Have a simple way to maintain common files across different repos
    in sync through links and automatically diff-ing files

  \item Support for both local and remote development using IDEs (e.g.,
    PyCharm, Visual Studio Code)
\end{itemize}

% ----------------------------------------------------------------------------
\subsubsection{Python package management}

\begin{itemize}
  \item Carefully manage and control dependencies using Python managers
    (such as \texttt{poetry}) and virtual environments

  \item Code and containers can be versioned and kept in sync automatically
    since a certain version of the code can require a certain version of
    the container to run properly

    \begin{itemize}
      \item Code is versioned through Git

      \item Each container has a \texttt{changelog.txt} that contains
        the current version and the history
    \end{itemize}
\end{itemize}

% ----------------------------------------------------------------------------
\subsubsection{Testing}

\begin{itemize}
  \item Run end-to-end tests using \texttt{pytest} by automatically
    discover tests based on dependencies and test lists, supporting the
    dependencies needed by different directories

  \item Native support for both children-containers (i.e., Docker-in-Docker)
    and sibling containers
\end{itemize}

% ----------------------------------------------------------------------------
\subsubsection{DevOps functionalities}

\begin{itemize}
  \item Support automatically different stages for container development

    \begin{itemize}
      \item E.g., \texttt{test} / \texttt{local}, \texttt{dev},
        \texttt{prod}
    \end{itemize}

  \item Standardize ways of building, testing, retrieving, and
    deploying containers

  \item Ensure alignment between development environment, deployment, and
    CI/CD systems (e.g., GitHub Actions)

  \item Bootstrap the development system using a ``thin environment'',
    which has the minimum number of dependencies to allow development
    and deployment in exactly the same way in different setups (e.g., server,
    personal laptop, CI/CD)

  \item Manage dependencies in a way that is uniform across platforms and
    OSes, using Docker containers

  \item Separate the need to:

    \begin{itemize}
      \item Build and deploy containers (by devops)

      \item Use containers to develop and test (by developers)
    \end{itemize}

  \item Built-in support for multi-architecture builds (e.g, for Intel
    \texttt{x86} and Arm) across different OSes supporting containers (e.g.,
    Linux, MacOS, Windows Subsystem for Linux WSL)

  \item Support for developing, testing, and deploying multi-container
    applications
\end{itemize}

% ----------------------------------------------------------------------------
\subsubsection{Runnable directory}

The core concept of the proposed approach is a \textbf{runnable
directory} --- a self-contained, independently executable directory with
code, equipped with a dedicated DevOps setup. A repository is thus a special
case of a runnable directory. Developers typically work within a
single runnable directory for a given application, enabling them to test
and deploy code without affecting other parts of the codebase.

A runnable directory can contain other runnable directories as subdirectories.
For example, Figure 1 depicts three runnable directories: \texttt{A},
\texttt{B}, and \texttt{C}. Here, \texttt{A} and \texttt{C} are repositories,
with \texttt{C} incorporated into \texttt{A} as a submodule, while \texttt{B}
is a subdirectory within \texttt{A}. This setup provides the same
accessibility as if all the code were hosted in a single monorepo. Note
that each of \texttt{A}, \texttt{B}, and \texttt{C} has its own DevOps
pipeline --- a key feature of our approach, which is discussed further
in Section 3.2.

% rendered_images:begin
% ```mermaid
% graph RL
%     subgraph A [Runnable Dir A]
%         direction TB
%         subgraph C1 [Runnable Dir C]
%             DevOpsC1[DevOps C]
%         end
%         subgraph B [Runnable Dir B]
%             DevOpsB[DevOps B]
%         end
%         DevOpsA[DevOps A]
%     end
%     subgraph C [Runnable Dir C]
%         DevOpsC[DevOps C]
%     end
% 
%     C -->|Submodule| C1
% 
%     style A fill:#FFF3CD
%     style C fill:#FFF3CD,stroke:#9E9D24
% ```
% label=fig:Runnable_Dir_Arch
% caption=Sample architecture of Causify's runnable directories.
% rendered_images:end
% render_images:begin
\begin{figure}[H]
  \includegraphics[width=\linewidth]{part2.runnable_directory.tex.figs/part2.runnable_directory.1.png}
  \caption{Sample architecture of Causify's runnable directories.}
  \label{fig:Runnable_Dir_Arch}
\end{figure}
% render_images:end

% ----------------------------------------------------------------------------
\subsubsection{Docker}

Docker is the backbone of our containerized development environment. Every
runnable directory contains Dockerfiles that allow it to build and run
its own Docker containers, which include the code, its dependencies,
and the runtime system.

This Docker-based approach addresses two important challenges. First,
it ensures consistency by isolating the application from variations in
the host operating system or underlying infrastructure. Second, a
specific package (or package version) can be added to the container of
a particular runnable directory without affecting other parts of the
codebase. This prevents ``bloating'' the environment with packages required
by all applications --- a common issue in monorepos --- while also effectively
mitigating the risk of conflicting dependencies, which can arise in a
multi-repo setup.

Our approach supports multiple stages for container release:

\begin{itemize}
  \item Local: used to work on updates to the container; accessible only
    to the developer who built it.

  \item Development: used by all team members in day-to-day
    development of new features.

  \item Production: used to run the system by end users.
\end{itemize}

This multi-stage workflow enables seamless progression from testing to
system deployment.

It is also possible to run a container within another container's environment
in a Docker-in-Docker setup. In this case, children containers are started
directly inside a parent container, allowing nested workflows or
builds. Alternatively, sibling containers can run side by side and
share resources such as the host's Docker daemon, enabling inter-container
communication and orchestration.

% rendered_images:begin
% ```mermaid
% graph TD
%     host[Host]
%     docker_engine[Docker Engine]
%     subgraph sibling_container["Sibling Containers"]
%         container_1[Container 1]
%         container_2[Container 2]
%     end
%     subgraph children_container["Children Containers"]
%         container_1a[Container 1a]
%         container_1b[Container 1b]
%     end
%     host --> docker_engine
%     docker_engine --> container_1
%     docker_engine --> container_2
%     container_1 --> container_1a
%     container_1 --> container_1b
% 
%     style sibling_container fill:#FFF3CD,stroke:#9E9D24
%     style children_container fill:#FFF3CD,stroke:#9E9D24
% ```
% caption=Docker container flow.
% label=fig:Docker_Flow
% rendered_images:end
% render_images:begin
\begin{figure}[H]
  \includegraphics[width=\linewidth]{part2.runnable_directory.tex.figs/part2.runnable_directory.2.png}
  \caption{Docker container flow.}
  \label{fig:Docker_Flow}
\end{figure}
% render_images:end

% ----------------------------------------------------------------------------
\subsubsection{What is needed}

An ideal strategy would combine the best of both worlds:

\begin{itemize}
  \item The modularity of multi-repos, to keep the codebase scalable
    and simplify day-to-day development processes.

  \item The environment consistency of monorepos, to avoid
    synchronization issues and prevent errors that arise from
    executing code in misaligned environments.
\end{itemize}

Both are achieved through the hybrid approach proposed in this paper,
which will be discussed in Section 3.

% ----------------------------------------------------------------------------
\subsubsection{Thin environment}

To bootstrap development workflows, we use a thin client that installs
a minimal set of essential dependencies, such as Docker and invoke, in
a lightweight virtual environment. A single thin environment is shared
across all runnable directories which minimizes setup overhead (see Figure
3). This environment contains everything that is needed to start development
containers, which are in turn specific to each runnable directory. With
this approach, we ensure that development and deployment remain
consistent across different systems (e.g., server, personal laptop, CI/CD).

% rendered_images:begin
% ```mermaid
% graph RL
%   thin_env[thin environment]
%   subgraph A [Runnable Dir A]
%     direction TB
%         B[Runnable Dir B]
%         C1[Runnable Dir C]
%     end
%     subgraph C [Runnable Dir C]
%     end
% 
%   C -->|Submodule| C1
%   A -.-> thin_env
%   B -.-> thin_env
%   C1 -.-> thin_env
%   C -.-> thin_env
% 
%   style A fill:#FFF3CD
%   style C fill:#FFF3CD,stroke:#9E9D24
% ```
% caption=Thin environment shared across multiple runnable directories.
% rendered_images:end
% render_images:begin
\begin{figure}[H]
  \includegraphics[width=\linewidth]{part2.runnable_directory.tex.figs/part2.runnable_directory.3.png}
  \caption{Thin environment shared across multiple runnable directories.}
\end{figure}
% render_images:end

% ----------------------------------------------------------------------------
\subsubsection{Submodule of ``helpers''}

All Causify repositories include a dedicated ``helpers'' repository as
a submodule. This repository contains common utilities and development
toolchains, such as the thin environment, Linter, Docker, and invoke
workflows. By centralizing these resources, we eliminate code
duplication and ensure that all teams, regardless of the project, use
the same tools and procedures.

Additionally, it hosts symbolic link targets for files that must technically
reside in each repository but are identical across all of them (e.g.,
license and certain configuration files). Manually keeping them in
sync can be difficult and error-prone over time. In our approach,
these files are stored exclusively in ``helpers'', and all other repositories
utilize read-only symbolic links pointing to them. This way, we avoid file
duplication and reduce the risk of introducing accidental discrepancies.

% rendered_images:begin
% ```mermaid
% graph RL
%     subgraph A [Runnable Dir A]
%         direction TB
%         B[Runnable Dir B]
%         H1[Helpers]
%     end
%     subgraph H [Helpers]
%     end
% 
%     H -->|Submodule| H1
% 
%     style A fill:#FFF3CD
%     style H fill:#FFF3CD,stroke:#9E9D24
% ```
% caption="Helpers" submodule integrated into a repository.
% rendered_images:end
% render_images:begin
\begin{figure}[H]
  \includegraphics[width=\linewidth]{part2.runnable_directory.tex.figs/part2.runnable_directory.4.png}
  \caption{"Helpers" submodule integrated into a repository.}
\end{figure}
% render_images:end

\paragraph{3.4.1. Git hooks}

Our ``helpers'' submodule includes a set of Git hooks used to enforce policies
across our development process, including Git workflow rules, coding standards,
security and compliance, and other quality checks. These hooks are
installed by default when the user activates the thin environment. They
perform essential checks such as verifying the branch, author
information, file size limits, forbidden words, Python file compilation,
and potential secret leaks\ldots etc.

% ----------------------------------------------------------------------------
\subsubsection{Executing tests}

Our system supports robust testing workflows that leverage the
containerized environment for comprehensive code validation. Tests are
executed inside Docker containers to ensure consistency across development
and production environments, preventing discrepancies caused by
variations in host system configurations. In the case of nested runnable
directories, tests are executed recursively within each directory's
corresponding container, which is automatically identified (see Figure
5). As a result, the entire test suite can be run with a single command,
while still allowing tests in subdirectories to use dependencies that
may not be compatible with the parent directory's environment.

% rendered_images:begin
% ```mermaid
% graph LR
%     start((start))
%     start --> A
%     subgraph A[Runnable Dir A]
%         direction LR
%         pytest_1((pytest))
%         B[Runnable Dir B / Container B]
%         C[Runnable Dir C / Container C]
%         dirA1[dir1 / Container A]
%         dirA2[dir2 / Container A]
%         dirA11[dir1.1 / Container A]
%         dirA12[dir1.2 / Container A]
%         pytest_1 --> B
%         pytest_1 --> C
%         pytest_1 --> dirA1
%         pytest_1 --> dirA2
%         dirA1 --> dirA11
%         dirA1 --> dirA12
%     end
% 
% style A fill:#FFF3CD,stroke:#9E9D24
% style B font-size:15px
% style C font-size:15px
% caption=Recursive test execution in dedicated containers.
% 
% ----------------------------------------------------------------------------
% \subsubsection{Dockerized executables}
% 
% Sometimes, installing a toolchain inside a development container is
% not justified especially when it is large or used only occasionally. In
% such cases, we use \emph{dockerized executables} where the tool and its
% dependencies are packaged into a slim Docker image, and the command is
% executed inside a short-lived container. We launch a short-lived container
% with the repository mounted, run the tool, then remove the container;
% the versioned image stays cached for reuse, avoiding repeated setup.
% 
% This approach prevents the development environment from being bloated with
% rarely used dependencies while preserving reproducibility and version control
% of the toolchain. When needed (e.g., during test execution), a dockerized
% executable can also run inside another Docker container using either the
% children or sibling container approach, as discussed in Section~3.2.
% 
% ============================================================================
% \subsection{Discussion}
% 
% Causify's approach presents a strong alternative to existing code organization
% solutions, offering scalability and efficiency for both small and
% large systems.
% 
% The proposed modular architecture is centered around runnable
% directories, which operate as independent units with their own build and
% release lifecycles. This design bypasses the bottlenecks common in
% large monorepos, where centralized workflows can slow down CI/CD processes
% unless specialized tools like Buck or Bazel are used. By leveraging
% Docker containers, we ensure consistent application behavior across development,
% testing, and production environments, avoiding problems caused by
% system configuration discrepancies. Dependencies are isolated within each
% directory's dedicated container, reducing the risks of issues that
% tight coupling or package incompatibility might create in a monorepo
% or a multi-repo setup.
% 
% Unlike multi-repos, runnable directories can utilize shared utilities from
% ``helper'' submodules, eliminating code duplication and promoting
% consistent workflows across projects. They can even reside under a unified
% repository structure which simplifies codebase management and reduces the
% overhead of maintaining multiple repositories. With support for
% recursive test execution spanning all components, runnable directories
% allow for end-to-end validation of the whole codebase through a single
% command, removing the need for testing each repository separately.
% 
% There are, however, several challenges that might arise in the
% adoption of our approach. Teams that are unfamiliar with containerized
% environments may need time and training to effectively transition to the
% new workflows. The reliance on Docker may introduce additional
% resource demands, particularly when running multiple containers concurrently
% on development machines. This would require further optimization, possibly
% aided by customized tooling. These adjustments, while ultimately beneficial,
% can add complexity to the system's rollout and necessitate ongoing
% maintenance to ensure seamless integration with existing CI/CD
% pipelines and development practices.
% 
% ============================================================================
% \subsection{Future directions}
% 
% Looking ahead, there are several areas where the proposed approach can
% be improved. One direction is the implementation of dependency-aware caching
% to ensure that only the necessary components are rebuilt or retested when
% changes are made. This would reduce the time spent on development
% tasks, making the overall process more efficient. Further optimization
% could involve designing our CI/CD pipelines to execute builds, tests,
% and deployments for multiple runnable directories in parallel, which would
% allow us to take full advantage of available compute resources.
% 
% Additional measures can also be taken to enhance security. Integrating
% automated container image scanning and validation before deployment
% would help guarantee compliance with organizational policies and
% prevent vulnerabilities from entering production environments. In addition,
% fine-grained access controls could be introduced for runnable
% directories in order to safeguard sensitive parts of the codebase. These
% steps will bolster both the security and efficiency of our workflows as
% the projects continue to scale.
% 
% 

  % ############################################################################
  % ############################################################################
\section{Dockerized Executables}

% Related documentation files:
% - helpers_root/docs/tools/docker/all.dockerized_flow.explanation.md
%   Core explanation of dockerized executable pattern
% - helpers_root/docs/tools/dev_system/all.docker_optimizer_container.how_to_guide.md
%   Specific example of dockerized executable for optimization
% - helpers_root/docs/tools/documentation_toolchain/all.latex_toolchain.how_to_guide.md
%   LaTeX as a dockerized executable
% - helpers_root/docs/tools/documentation_toolchain/all.render_images.explanation.md
%   Image rendering tools as dockerized executables
% - helpers_root/docs/tools/documentation_toolchain/all.notes_to_pdf.explanation.md
%   PDF generation as dockerized executable
% - helpers_root/docs/tools/all.llm_transform.how_to_guide.md
%   LLM-powered transformations as dockerized executables

% ============================================================================
\subsection{Motivation}

Software development teams frequently encounter the challenge of managing
diverse toolchains across heterogeneous environments. Installing and
maintaining developer utilities often leads to version conflicts, platform-specific
incompatibilities, and inconsistent behavior across development
machines, CI/CD pipelines, and production systems. These issues compound
when teams work across multiple operating systems (macOS, Linux, WSL) and
deployment contexts (local development, containerized environments,
cloud infrastructure).

Traditional approaches to tool management present several limitations.
Installing tools directly on host systems creates fragile dependencies
that differ across platforms. Packaging all utilities into a single monolithic
development container leads to bloated images, lengthy build times, and
dependency conflicts between tools. Version management becomes complex
when different projects require different tool versions, and onboarding
new team members involves extensive setup documentation that quickly becomes
outdated.

To address these challenges, we developed \textbf{Dockerized
Executables}---a pattern that packages each development tool in its own
slim, purpose-built container image. Instead of installing tools locally
or embedding them in a single development container, each utility runs
on-demand in an isolated environment with all dependencies pre-configured.
This approach provides clean separation between tools, enables
independent versioning and upgrades, reduces security surface area,
and ensures consistent behavior across all execution contexts.

% ============================================================================
\subsection{Summary}

Dockerized Executables solve a persistent problem in software
development: the challenge of maintaining consistent, reproducible
tooling across diverse environments. By packaging each tool in an isolated
container with a transparent wrapper, we eliminate installation
complexity, ensure consistent behavior, and maintain clean development
environments. While this approach requires Docker as a dependency, the
benefits in reproducibility, onboarding speed, and operational
simplicity justify this requirement. The pattern has proven effective across
diverse tooling needs, from document processing to diagram generation
to PDF compilation, and has become a foundational component of our development
infrastructure.

% ============================================================================
\subsection{Architecture and Design}

A Dockerized Executable consists of two primary components: a minimal
container image containing the tool and its dependencies, and a thin wrapper
script that handles execution logistics. From the developer's
perspective, invoking a Dockerized Executable is indistinguishable from
running a locally installed tool---the wrapper transparently manages containerization
details.

% ----------------------------------------------------------------------------
\subsubsection{Container image structure}

Each tool is packaged in a dedicated Docker image built on a minimal
base (e.g., \texttt{debian:stable-slim}). The image includes only the
dependencies required for that specific utility, keeping the image size
small and reducing attack surface. Images are version-pinned and
tagged to enable controlled upgrades and rollbacks. The container runs
with a non-root user by default and uses a standardized working directory
(\texttt{/src}) where the repository is mounted.

% ----------------------------------------------------------------------------
\subsubsection{Wrapper script responsibilities}

The wrapper script bridges between the user's environment and the containerized
tool. Key responsibilities include:

\begin{itemize}
  \item \textbf{Repository discovery}: Locates the repository root to establish
    a consistent mount point regardless of the user's current working directory.

  \item \textbf{Path translation}: Converts file paths between host/dev-container
    coordinates and container-internal paths (e.g., \texttt{docs/file.md}
    → \texttt{/src/docs/file.md}).

  \item \textbf{Container lifecycle management}: Pulls or builds the image
    if needed (leveraging Docker's caching), launches the container
    with appropriate mounts and permissions, and removes it after execution.

  \item \textbf{Permission preservation}: Runs the container with the caller's
    UID:GID to ensure generated files have correct ownership.

  \item \textbf{Environment isolation}: Whitelists and forwards only necessary
    environment variables, preventing unintended dependencies on host
    configuration.

  \item \textbf{Output handling}: Streams stdout/stderr in real-time and
    propagates the tool's exit code for proper CI/CD integration.
\end{itemize}

% ----------------------------------------------------------------------------
\subsubsection{Repository root as anchor point}

Using the repository root as the mount anchor is critical for
consistency. This design choice ensures:

\begin{itemize}
  \item Commands work identically regardless of the user's current subdirectory

  \item Paths remain stable across local development, dev containers, and
    CI environments

  \item The same command works in both children and sibling container patterns

  \item Minimal, predictable mounts prevent exposing unrelated host
    directories
\end{itemize}

% ============================================================================
\subsection{Execution Flow}

When a developer invokes a Dockerized Executable:

\begin{enumerate}
  \item The wrapper script is called with the tool's arguments

  \item The script discovers the repository root and normalizes input paths

  \item It ensures the tool's Docker image is available (pulling/building
    if necessary, leveraging cache otherwise)

  \item The repository is mounted into a new container at \texttt{/src}

  \item Paths are rewritten to container coordinates

  \item The tool executes inside the container with the provided arguments

  \item Outputs are written back to the mounted repository

  \item The container exits, and its exit code is propagated to the
    caller
\end{enumerate}

The entire process is transparent to the user, who sees only the tool's
output and return code.

% ============================================================================
\subsection{Container Execution Patterns}

Dockerized Executables must operate correctly whether invoked from the
host system or from within development containers. This requirement introduces
complexity in mounting and Docker daemon access.

% ----------------------------------------------------------------------------
\subsubsection{Children containers (Docker-in-Docker)}

In this pattern, a development container runs its own Docker daemon and
spawns tool containers as children. While flexible, this approach requires
elevated privileges (typically \texttt{--privileged}) and introduces
security concerns. It also adds overhead from running nested Docker daemons.

% ----------------------------------------------------------------------------
\subsubsection{Sibling containers (preferred)}

Our preferred approach mounts the host's Docker socket into the
development container, allowing it to launch sibling containers via the
host's Docker engine. The development container and tool container run
as peers, both managed by the host daemon. This pattern is more secure,
more efficient, and avoids privilege escalation. However, it requires
careful path management: mounts must reference host-visible paths, not
dev-container-internal paths. The wrapper script handles this translation
automatically.

% ============================================================================
\subsection{Practical Examples}

We have successfully Dockerized numerous development utilities:

\begin{itemize}
  \item \textbf{Document formatting}: Prettier for consistent code and
    Markdown formatting across the codebase

  \item \textbf{Document conversion}: Tools for converting between formats
    (e.g., DOCX to Markdown) without installing converters locally

  \item \textbf{Diagram rendering}: Mermaid, Graphviz, and TikZ renderers
    for generating diagrams from source specifications

  \item \textbf{LaTeX compilation}: Complete LaTeX toolchain for building
    PDFs without installing LaTeX distributions on development machines

  \item \textbf{LLM-powered transforms}: Utilities leveraging language
    models for structured content transformations
\end{itemize}

Each tool follows the same pattern: a single Python wrapper script, a minimal
Dockerfile, and standardized invocation semantics.

% ----------------------------------------------------------------------------
\subsubsection{Example invocation}

Converting a document requires only:

\begin{verbatim}
python dev_scripts_helpers/documentation/convert_docx_to_markdown.py \
-i docs/weekly_update.docx \
-o docs/weekly_update.md
\end{verbatim}

The wrapper handles image management, path translation, and execution
automatically. The result is identical to a local installation, but
with consistent behavior across all platforms.

% ============================================================================
\subsection{Benefits and Trade-offs}

The Dockerized Executable pattern provides several advantages:

\begin{itemize}
  \item \textbf{Reproducibility}: Identical tool behavior across all environments
    (macOS, Linux, WSL, CI)

  \item \textbf{Rapid onboarding}: Docker is the only prerequisite; no
    lengthy installation guides

  \item \textbf{Clean environments}: Development containers remain lean,
    containing only essential dependencies

  \item \textbf{Independent versioning}: Each tool can be upgraded or rolled
    back without affecting others

  \item \textbf{Reduced conflicts}: Isolated dependencies prevent version
    collisions between tools

  \item \textbf{Smaller security surface}: Purpose-built images contain
    only necessary components
\end{itemize}

However, this approach introduces some costs:

\begin{itemize}
  \item \textbf{Docker dependency}: The system requires Docker to be available
    and properly configured

  \item \textbf{Image storage}: Each tool requires disk space for its image,
    though images are typically small (tens to hundreds of MB)

  \item \textbf{Initial pull latency}: First-time execution requires pulling
    or building the image, though subsequent runs leverage Docker's
    cache

  \item \textbf{Slight overhead}: Container startup adds minimal latency
    (typically under 1 second) compared to native execution
\end{itemize}

For tools used frequently in development workflows, these trade-offs are
favorable. The reproducibility and consistency benefits far outweigh
the minimal overhead.

% ============================================================================
\subsection{Comparison with Language-Specific Solutions}

The Dockerized Executable pattern addresses a broader problem space than
language-specific dependency managers. Tools like \texttt{uv} (for
Python) provide fast, isolated dependency management within their
language ecosystem. They excel at version pinning and environment isolation
for language packages but are limited to that language's tooling.

Dockerized Executables complement such tools rather than replacing them.
They provide consistent execution for:

\begin{itemize}
  \item Tools written in languages not used in the main project

  \item System-level utilities with complex native dependencies

  \item Tools requiring specific OS packages or libraries

  \item Cross-language workflows requiring coordinated toolchains
\end{itemize}

Both approaches share the goal of reproducible, isolated tool
execution. The choice depends on whether the tool fits within a language
ecosystem's package manager or requires the broader isolation that containers
provide.

% ============================================================================
\subsection{Implementation Guidelines}

Building a Dockerized Executable follows a consistent pattern:

\begin{enumerate}
  \item \textbf{Create a minimal image}: Start with a small base image
    (\texttt{debian:stable-slim}, \texttt{alpine}, or language-specific
    slim variants). Install only necessary dependencies. Pin versions explicitly.
    Use a non-root user. Set \texttt{WORKDIR /src}.

  \item \textbf{Build the wrapper script}: Implement repository root discovery.
    Mount the repo at \texttt{/src}. Translate paths from host to
    container coordinates. Forward necessary arguments and environment
    variables. Run with caller's UID:GID. Stream logs and propagate exit
    codes.

  \item \textbf{Choose execution pattern}: Prefer sibling containers for
    security and efficiency. Handle path translation correctly for
    sibling pattern. Document any special requirements.

  \item \textbf{Test thoroughly}: Invoke from repository root and subdirectories.
    Test from host and dev containers. Verify path handling and file ownership.
    Assert on outputs and exit codes.

  \item \textbf{Document usage}: Provide clear examples. Document
    required environment variables if any. Note any platform-specific considerations.
\end{enumerate}

% ============================================================================
\subsection{Integration with Development Workflows}

Dockerized Executables integrate seamlessly into existing workflows:

\begin{itemize}
  \item \textbf{Local development}: Developers invoke tools via wrapper
    scripts, receiving immediate feedback

  \item \textbf{Pre-commit hooks}: Git hooks can invoke Dockerized Executables
    for formatting or validation

  \item \textbf{CI/CD pipelines}: GitHub Actions and other CI systems invoke
    the same wrapper scripts, ensuring consistency

  \item \textbf{Automated workflows}: Tools can be orchestrated via make,
    invoke, or other task runners

  \item \textbf{Interactive use}: Developers can use tools ad-hoc during
    development without environment concerns
\end{itemize}

This consistency across contexts eliminates "works on my machine" issues
and simplifies troubleshooting.

% ############################################################################
\section{Docker container release flow}

% Related documentation files:
% - docs/infra/ck.development_stages.explanation.md
%   Explains dev/prod/local container stages
% - docs/infra/ck.elastic_container_service.explanation.md
%   ECS setup for container deployment
% - docs/infra/ck.ecs_flows.reference.md
%   Reference for ECS container flows
% - docs/infra/ck.aws_ecr.explanation.md
%   AWS ECR for container registry management
% - docs/infra/ck.airflow_dags_release.how_to_guide.md
%   Airflow DAG release process coordinated with containers
% - helpers_root/docs/tools/dev_system/all.devops_docker_auto_release.explanation.md
%   Automated Docker container release workflows
% - helpers_root/docs/tools/dev_system/all.devops_docker.how_to_guide.md
%   Manual Docker release procedures
% - helpers_root/docs/tools/dev_system/all.devops_docker.reference.md
%   Complete reference for container release system
% - docs/build/ck.gh_workflows.explanation.md
%   GitHub workflows for building containers
% - docs/build/ck.gh_workflows.reference.md
%   Reference for all build workflows
% - docs/bwd/ck.bwd_app_dev_image_release.how_to_guide.md
%   Development image release process
% - docs/bwd/ck.bwd_app_backend_release.how_to_guide.md
%   Backend production container release
% - docs/trading_ops/ck.release.explanation.md
%   Release process including container updates
% - docs/trading_ops/ck.release.how_to_guide.md
%   Step-by-step release guide
% - docs/trading_ops/ck.dags_release.how_to_guide.md
%   DAG release synchronized with container versions
% - docs/all.dev_system_release.how_to_guide.md
%   Overall development system release process

% ============================================================================
\subsection{Motivation and Architecture}
Modern software systems require consistent, reproducible deployment environments
across development, staging, and production. Traditional deployment
approaches struggle with environmental drift, dependency conflicts,
and versioning complexity. Docker containers provide isolation and reproducibility,
but managing container lifecycles—building, testing, versioning, and deploying—introduces
operational overhead. Our Docker container release flow addresses
these challenges through automated workflows that enforce quality
gates, maintain version consistency, and provide staged rollout
capabilities. The system balances developer velocity with operational safety,
enabling rapid iteration while preventing production incidents. The architecture
distinguishes between multiple container types and deployment stages:
\begin{itemize}
  \item \textbf{Development images} (\texttt{dev}): Used for day-to-day
    development, containing all development dependencies, debugging
    tools, and test frameworks. Updated frequently (daily or on-demand).

  \item \textbf{Production images} (\texttt{prod}): Optimized for
    deployment, containing only runtime dependencies with minimal attack
    surface. Updated on a controlled schedule (bi-weekly or when critical
    fixes are required).

  \item \textbf{Local images} (\texttt{local}): Developer-specific
    builds for testing container changes before promoting to dev. Never
    pushed to registries.
\end{itemize}
Each image stage has corresponding task definitions in AWS ECS (Elastic
Container Service) for orchestrated deployments. The entire flow is managed
through GitHub Actions workflows that run on push, schedule, or manual
dispatch.
% ============================================================================
\subsection{Container Build Pipeline}
The container build pipeline transforms source code and dependencies
into versioned, tested Docker images ready for deployment. The process
is automated through GitHub Actions workflows that run in response to
code changes or manual triggers.
% ----------------------------------------------------------------------------
\subsubsection{Development Image Workflow}
Development images are built and released through the \texttt{build\_image.dev.yml}
workflow. This workflow:
\begin{enumerate}
  \item Triggers on:
    \begin{itemize}
      \item Push to master branch

      \item Manual workflow dispatch with optional version override

      \item Scheduled runs (daily for routine updates)
    \end{itemize}

  \item Checks out the repository at the specified commit or branch

  \item Reads the current version from \texttt{changelog.txt}, which maintains
    a human-readable history of container changes with semantic versioning
    (major.minor.patch)

  \item Builds the image using the Dockerfile in the repository root,
    installing development dependencies, test frameworks, and
    debugging tools

  \item Tags the image with:
    \begin{itemize}
      \item Version tag (e.g., \texttt{1.2.3})

      \item \texttt{dev} tag pointing to the latest development
        version

      \item Commit SHA for precise traceability
    \end{itemize}

  \item Pushes the image to Amazon ECR (Elastic Container Registry)

  \item Updates environment variables and configuration to reference the
    new image version

  \item Optionally triggers downstream workflows to update dependent systems
\end{enumerate}
The development image build is relatively fast (typically 10-20 minutes)
because it leverages Docker layer caching and does not require
extensive optimization passes.
% ----------------------------------------------------------------------------
\subsubsection{Production Image Workflow}
Production images follow a similar but more rigorous process through \texttt{build\_image.cmamp.yml}
(or equivalent production workflow). Key differences include:
\begin{itemize}
  \item \textbf{Stricter trigger conditions}: Only builds on manual dispatch
    or scheduled releases, never on every push. This prevents immature
    code from reaching production workflows.

  \item \textbf{Multi-stage builds}: Uses Docker multi-stage builds to
    minimize final image size. Build dependencies are discarded, and only
    runtime requirements remain in the final image.

  \item \textbf{Security scanning}: Integrates vulnerability scanning (e.g.,
    Trivy, AWS ECR scanning) to detect known CVEs in base images and
    dependencies before release.

  \item \textbf{Optimization passes}: Removes unnecessary files,
    compresses layers, and ensures minimal attack surface

  \item \textbf{Comprehensive testing}: Runs full test suites (fast, slow,
    and sometimes superslow) to validate the production build before release

  \item \textbf{Version locking}: Production images are immutable once
    tagged. Rollbacks occur by redeploying a previous version, not by modifying
    existing images.
\end{itemize}
Production image builds are slower (30-60 minutes) due to these
additional steps, but the investment ensures reliability and security.
% ----------------------------------------------------------------------------
\subsubsection{Version Management}
Container versions are managed through \texttt{changelog.txt} files that
serve as the single source of truth for the current version. Each
entry includes:
\begin{itemize}
  \item Version number (semantic versioning: major.minor.patch)

  \item Release date

  \item Summary of changes (features, fixes, dependency updates)

  \item Author information
\end{itemize}
When a developer needs to release a new version, they:
\begin{enumerate}
  \item Update \texttt{changelog.txt} with the new version and
    description

  \item Commit the change to the repository

  \item Trigger the build workflow (automatically on push to master or
    manually via workflow dispatch)

  \item The workflow reads the version from \texttt{changelog.txt} and
    tags the image accordingly
\end{enumerate}
This approach ensures version changes are tracked in version control and
reviewed during code review, preventing accidental or undocumented
version bumps.
% ============================================================================
\subsection{Task Definition Management with ECS}
AWS Elastic Container Service (ECS) orchestrates containerized applications
using task definitions that specify which Docker images to run,
resource allocations, networking configuration, and environment
variables. Our system manages separate task definitions for preproduction
and production environments, enabling staged rollouts and safe testing
of changes.

% ----------------------------------------------------------------------------
\subsubsection{Preproduction Task Definition Release}
The \texttt{release\_new\_ecs\_preprod\_task\_definition.yml} workflow
updates the preproduction ECS task definition. This workflow:
\begin{enumerate}
  \item Triggers on manual dispatch, allowing controlled releases
    after image builds complete

  \item Fetches the current task definition from AWS ECS

  \item Updates the image reference to point to the newly built development
    image (specified by version or tag)

  \item Modifies environment variables if needed (e.g., database
    connection strings for preproduction)

  \item Registers the new task definition revision with ECS

  \item Optionally forces a new deployment of running services using this
    task definition, causing ECS to restart tasks with the new image

  \item Logs the task definition ARN and revision number for traceability
\end{enumerate}
Preproduction task definitions use development images and connect to
isolated resources (databases, message queues) to avoid impacting production
systems. This environment serves as the final validation stage before
production deployment.

% ----------------------------------------------------------------------------
\subsubsection{Production Task Definition Release}
The \texttt{release\_new\_ecs\_prod\_task\_definition.yml} workflow
follows an identical structure but with additional safeguards:
\begin{itemize}
  \item Requires explicit manual approval (via GitHub environment
    protection rules) before proceeding

  \item Only references production-tagged images that have passed all quality
    gates

  \item Logs all changes to an audit trail for compliance and debugging

  \item Supports blue-green deployments and canary releases through ECS
    service configuration

  \item Monitors deployment progress and can automatically roll back if
    health checks fail
\end{itemize}
Production releases typically occur bi-weekly on a fixed schedule unless
critical patches are required. This cadence balances feature velocity
with operational stability.

% ============================================================================
\subsection{Airflow DAG Release Process}
Airflow DAGs (Directed Acyclic Graphs) define workflows for data processing,
model training, and automated trading operations. DAG code changes
often require updated Docker images containing new dependencies or
business logic. The DAG release process coordinates code changes with container
updates.

% ----------------------------------------------------------------------------
\subsubsection{Release Workflow for Trading DAGs}
The release procedure for Airflow DAGs follows a structured, multi-step
process:
\begin{enumerate}
  \item \textbf{Create and test a DAG}: When updating DAG logic,
    create a test DAG with a task-specific name (e.g., \texttt{test.tokyo.shadow\_trading.C11a.config1\_CmTask8080}).
    This naming convention embeds the task number for traceability.

  \item \textbf{Validate with preproduction task definition}: Execute the
    test DAG using the existing preproduction ECS task definition. Do
    not update the image yet—this isolates DAG logic changes from
    environment changes. Run for a minimum duration (e.g., 30 minutes)
    to ensure stability.

  \item \textbf{Verify all tasks succeed}: Check that every task in the
    DAG completes successfully without errors or warnings. Review logs
    for unexpected behavior.

  \item \textbf{Propagate changes to production DAGs}: Once validated,
    apply the changes to all relevant production DAGs in \texttt{datapull/airflow/dags/trading/shadow\_trading/}.

  \item \textbf{Update shared utilities}: If new utility functions were
    added to \texttt{datapull/airflow/dags/airflow\_utils/trade\_exec/utils.py},
    copy them to the shared Airflow utilities directory (\texttt{/data/shared/airflow/dags/airflow\_utils/trade\_exec/utils.py})
    after validation.

  \item \textbf{Merge and document}: Open a pull request for review.
    Once approved, merge to master. Create a release issue titled ``Release
    shadow trading DAGs — \{date\}'' listing affected DAGs and links
    to merged PRs.

  \item \textbf{Image update if needed}: If the DAG changes require a new
    Docker image (e.g., new Python packages, updated business logic), note
    this in the release issue and coordinate with the production release
    schedule.
\end{enumerate}
This process ensures DAG changes are validated independently before
deployment and that dependencies between DAGs, images, and
infrastructure are clearly managed.

% ----------------------------------------------------------------------------
\subsubsection{Release Schedule}
Shadow trading DAGs follow a bi-weekly release cadence when image updates
are required. When only DAG logic changes (no new dependencies), releases
occur weekly. This schedule provides predictability for operations
teams while allowing rapid iteration on trading strategies.

% ============================================================================
\subsection{Feature Release Communication and Documentation}
Releasing new features or significant system changes requires more than
code deployment. Effective communication ensures users understand what
changed, why it matters, and how to use new capabilities. Our process integrates
documentation updates and stakeholder communication into the release workflow.

% ----------------------------------------------------------------------------
\subsubsection{Pre-Release Internal Testing}
Before broadly announcing a feature, we conduct internal testing with a
small group:
\begin{enumerate}
  \item \textbf{Draft communication}: Write a detailed description of the
    feature in a Google Doc, explaining what changed, why, expected
    user actions, and links to documentation or dashboards.

  \item \textbf{Team review}: Share the draft with the feature development
    team for feedback on technical accuracy and clarity.

  \item \textbf{Pilot rollout}: Send the communication to a limited group
    (feature team or power users) and ask for feedback on both the feature
    and the communication.

  \item \textbf{Iterate}: Incorporate feedback and resolve any issues discovered
    during the pilot.
\end{enumerate}
This approach catches problems early when they're easier to fix and
ensures communication is clear before reaching a wider audience.
% ----------------------------------------------------------------------------
\subsubsection{Updating Central Release Documentation}
After internal validation, update the master release tracking document
(\textit{Dev System Features} Google Doc) with:
\begin{itemize}
  \item Feature title

  \item Release date

  \item Summary of changes

  \item Links to detailed documentation, dashboards, or related issues
\end{itemize}
This central document serves as a historical record of system
evolution and a reference for onboarding new team members.
% ----------------------------------------------------------------------------
\subsubsection{Communication Template}
Release announcements follow a structured template:
\begin{verbatim}
Subject: [Feature Name] Now Available
Hi Team,
We've integrated [feature name] to [solve what problem].
What's Covered?
- Key capability 1
- Key capability 2
- Integration points and usage
Current State:
- Metrics or dashboards (with access links)
- Example visuals or screenshots
Next Steps:
- Recommended actions for users
- How to get help or provide feedback
For detailed documentation, see: [link]
Thanks,
[Team]
\end{verbatim}
This format ensures consistency, provides actionable information, and
reduces confusion about new features.
% ============================================================================
\subsection{Integration with CI/CD and Quality Gates}
The Docker release flow integrates tightly with our continuous integration
and deployment (CI/CD) pipelines. Quality gates at each stage prevent
broken code from progressing through the release pipeline.
% ----------------------------------------------------------------------------
\subsubsection{Automated Testing in Workflows}
Build workflows incorporate testing at multiple points:
\begin{itemize}
  \item \textbf{Pre-build validation}: Lint checks, code formatting, and
    static analysis run before building the image. If these fail, the build
    aborts immediately.

  \item \textbf{Post-build testing}: After the image is built, fast
    tests run inside the new container to validate that code executes
    correctly in the containerized environment. This catches environment-specific
    issues (e.g., missing dependencies, incorrect paths).

  \item \textbf{Integration testing}: For production builds,
    integration tests verify interactions between components (databases,
    APIs, message queues) using the new image.

  \item \textbf{Security scanning}: Vulnerability scanners analyze the
    image for known CVEs. High-severity vulnerabilities block the
    release until patched.
\end{itemize}
Each gate can block progression, ensuring only validated images reach deployment
stages.
% ----------------------------------------------------------------------------
\subsubsection{Rollback Capabilities}
Despite rigorous testing, issues sometimes arise in production. The
system supports rapid rollback:
\begin{itemize}
  \item \textbf{Immutable image tags}: Each version is permanently tagged
    in ECR. Rollback is as simple as updating the task definition to
    reference a previous version.

  \item \textbf{Git-tracked configuration}: Task definition changes are
    committed to the repository. To roll back, revert the commit and re-run
    the release workflow.

  \item \textbf{Automated health checks}: ECS monitors container health.
    If new tasks fail health checks repeatedly, ECS can automatically stop
    the rollout and maintain the previous version.

  \item \textbf{Blue-green deployments}: For critical services, we
    maintain parallel environments (blue and green). New versions deploy
    to the inactive environment, and traffic switches only after
    validation succeeds.
\end{itemize}
These mechanisms minimize downtime and reduce the blast radius of problematic
releases.
% ============================================================================
\subsection{Buildmeister Dashboard and Monitoring}
The Buildmeister Dashboard provides real-time visibility into the
state of all GitHub Actions workflows across all repositories. This centralized
view enables rapid identification of failures and simplifies daily operational
oversight. The dashboard displays:
\begin{itemize}
  \item \textbf{Workflow status}: Current state of all workflows (success,
    failed, in progress) across repositories

  \item \textbf{Links to runs}: Direct links to failing workflow runs for
    immediate investigation

  \item \textbf{Test reports}: Links to Allure reports showing test history,
    flaky tests, and failure trends

  \item \textbf{Coverage reports}: Latest code coverage metrics with drill-down
    by module

  \item \textbf{Pull request counts}: Number of open PRs per repository,
    indicating review backlog
\end{itemize}
The dashboard updates hourly and is published to an S3-hosted static site
accessible to all team members. Historical versions are retained for
analysis of workflow stability over time. This tool is central to the Buildmeister
role, which rotates weekly among team members. The Buildmeister monitors
the dashboard, triages failures, and ensures builds remain green.
% ============================================================================
\subsection{Summary}
Our Docker container release flow automates the journey from source
code to deployed containers while enforcing quality gates, maintaining
version consistency, and enabling safe rollouts. Separate workflows for
development and production images balance velocity with stability.
Integration with ECS task definitions provides orchestrated
deployments with rollback capabilities. The Airflow DAG release process
coordinates workflow updates with container changes. Structured
communication and documentation ensure stakeholders understand changes.
Finally, the Buildmeister Dashboard and automated monitoring provide
visibility into system health and enable rapid response to issues. This
system transforms container management from a manual, error-prone
process into a reliable, auditable pipeline. Developers can iterate
rapidly with development images while operations teams maintain
production stability through controlled releases. The architecture is designed
for growth: adding new services or repositories requires minimal
configuration, as all core workflows are reusable and parameterized.
Quality gates prevent broken code from reaching production, and comprehensive
monitoring ensures issues are detected and resolved quickly.


  % ############################################################################
  % ############################################################################
\section{AI-Optimized Development Infrastructure}

% ============================================================================
\subsection{Motivation}

The Causify development system is designed to enable AI-assisted development.
The same design principles that enable human developers to work efficiently
(such as discoverability, automation, standardization, and immediate feedback)
align remarkably well with the operational requirements and constraints of AI
coding agents. This section examines five key aspects of the system that create
an ``AI-optimized codebase'' and explains why these characteristics enable more
effective human-AI collaboration in software development.

% ============================================================================
\subsection{Self-Documenting, Executable Workflows}

The system's use of Python \texttt{invoke} tasks as a centralized task
registry fundamentally changes how developers (both human and AI) interact
with the codebase. Rather than requiring knowledge of complex command-line
invocations, Docker configurations, or repository-specific conventions,
all development operations are exposed through a discoverable, consistent
interface.

AI agents benefit from this design in several ways:

\begin{itemize}
  \item \textbf{Discoverability}: Agents can execute \texttt{invoke --list}
    to enumerate all available operations (testing, linting, deployment,
    Docker management, git operations) without prior knowledge of the codebase

  \item \textbf{Elimination of ambiguity}: Instead of constructing bash
    commands from documentation or examples that might be outdated, agents
    invoke standardized tasks with well-defined semantics

  \item \textbf{Cross-repository consistency}: The same commands work
    identically across all runnable directories, eliminating the need
    to learn repository-specific patterns

  \item \textbf{Encoded best practices}: The invoke tasks themselves
    represent the correct way to perform operations, including proper
    Docker container selection, pytest options, timeout configurations,
    and output handling
\end{itemize}

For example, rather than an AI agent attempting to construct the correct
Docker command with appropriate mounts, user permissions, and pytest
arguments, it simply executes \texttt{invoke run\_fast\_tests --coverage}.
This reduces the cognitive complexity the AI must manage and eliminates
entire categories of potential errors.

% =============================================================================
\subsection{Agent Instruction Manifests: Making the Repository Speak to AI}

In addition to executable workflows, the system embeds explicit, agent-targeted
guidance directly in the repository. The \texttt{csfy} repository and the nested
\texttt{helpers} module include a \texttt{CLAUDE.md} file that serves as an
\emph{agent instruction manifest}: a structured description of architecture,
module boundaries, and the canonical commands to run for testing, linting, and
container development.

This is not “documentation” in the traditional sense. It is a machine-consumable
operating manual that eliminates guesswork and prevents agents from inventing
unsafe or non-standard workflows. It also reduces prompt/context requirements:
instead of pasting a long “how to contribute” guide into an agent session, the
repository provides a stable, versioned source of truth.

\begin{lstlisting}[language=Text]
# Example excerpt (conceptual) from CLAUDE.md:
# - Architecture overview (components and ownership boundaries)
# - Canonical commands:
#   invoke run_fast_tests
#   invoke lint --modified
#   invoke docker_bash
\end{lstlisting}

Embedding instruction manifests yields two practical benefits:
(1) agents become immediately productive without an interactive onboarding loop,
and (2) maintainers can evolve agent behavior safely through code review
(change the manifest, not a hidden prompt).

% =============================================================================
\subsection{Machine-Readable Repository Contracts: \texttt{repo\_config.yaml}}

Runnable directories in \texttt{csfy} expose a machine-readable contract via
\texttt{repo\_config.yaml}. This file encodes the operational identity of the
directory: repository name, container image name, registries, and runnable-dir
behavior (including whether \texttt{helpers} is nested).

For AI agents, this is a major upgrade over prose documentation: it enables direct
queries like “what image name should I build?” or “which container registry should
I push to?” without searching or inference.

\begin{lstlisting}[language=YAML]
repo_info:
  repo_name: csfy
docker_info:
  docker_image_name: csfy
container_registry_info:
  ecr: <...>.dkr.ecr.<region>.amazonaws.com
  ghcr: ghcr.io/causify-ai
runnable_dir_info:
  use_helpers_as_nested_module: 1
  venv_tag: helpers
  dir_suffix: csfy
\end{lstlisting}

This contract makes automation safer. For example, a release task can derive the
target registry and image naming convention from \texttt{repo\_config.yaml} rather
than relying on environment-specific assumptions.

% =============================================================================
\subsection{Context-Bounded Editing: Runnable Directories as Prompt Boundaries}

AI agents operate under context constraints: they cannot reliably load or reason
about an entire monorepo at once. Runnable directories provide a natural
\emph{context boundary}. Each directory is intended to be self-contained: it has
its own Docker configuration, task entrypoints, and test lifecycle.

This turns an otherwise fuzzy question—“what subset of the repo matters?”—into a
concrete procedure:
\begin{itemize}
  \item focus the agent on a single runnable directory (or a small subtree),
  \item use its local \texttt{tasks.py}/invoke interface to build, test, and lint,
  \item treat other runnable directories as dependencies with explicit boundaries.
\end{itemize}

The result is practical: higher-quality agent edits with fewer accidental
cross-cutting changes and fewer “context hallucinations.”

% =============================================================================
\subsection{Guardrails and Provenance for Autonomous Changes}

Agentic development fails in predictable ways: style drift, accidental secret
inclusion, environment-specific behavior, and changes that “work locally” but fail
in CI. The system counters these with layered guardrails that are both
human-friendly and agent-friendly.

Key guardrails include:
\begin{itemize}
  \item \textbf{Formatting and lint automation}: pre-commit hooks (e.g., ruff) reduce
    style failures to deterministic fixes rather than review-time arguments.
  \item \textbf{Containerized execution}: tests run in containers, removing host
    variability as an explanation for failures.
  \item \textbf{Secret hygiene gates}: the workflow includes explicit secret-scanning
    gates (e.g., gitleaks in CI), ensuring “agent mistakes” do not become incidents.
  \item \textbf{Version synchronization}: environment and tooling versions are pinned
    and updated via controlled mechanisms (e.g., changelog/version workflows),
    preventing silent drift.
\end{itemize}

Equally important is \emph{provenance}: an agent should not merely propose code; it
should produce a change accompanied by evidence (test output, lint output, and a
clear summary of what changed and why). This supports the core practice-oriented
goal: make agent output reviewable and auditable.

% =============================================================================
\subsection{The Agentic Loop: Plan $\rightarrow$ Patch $\rightarrow$ Prove}

\begin{figure}[t]
  \centering
  \includegraphics[width=\linewidth]{papers/Causify_AIgentic_Development_System/figs/ai_optimized_development_infrastructure_flowchart.png}
  \caption{AI-optimized development infrastructure: instruction manifests (\texttt{CLAUDE.md}), machine-readable repo contracts (\texttt{repo\_config.yaml}), and runnable-directory boundaries feed a disciplined agentic loop (Plan$\rightarrow$Patch$\rightarrow$Prove$\rightarrow$Summarize), constrained by quality gates (pre-commit, lint, tests, secret scanning, version sync) to produce evidence-carrying, reviewable changes.}
  \label{fig:ai_optimized_dev_infra}
\end{figure}

The infrastructure enables a disciplined agent workflow:
\begin{enumerate}
  \item \textbf{Plan}: summarize intended change and affected runnable directory.
  \item \textbf{Patch}: implement a minimal diff that respects directory boundaries.
  \item \textbf{Prove}: run the smallest appropriate gates first, then expand:
    \begin{itemize}
      \item \texttt{invoke lint --modified}
      \item \texttt{invoke run\_fast\_tests} (or directory-scoped equivalent)
      \item optional: slow/superslow tests for risky changes
    \end{itemize}
  \item \textbf{Summarize}: provide a PR-ready explanation and evidence links/outputs.
\end{enumerate}

This loop is the practical bridge from “AI can generate code” to “AI can contribute
changes safely.” The novelty is not the existence of agents—it is the existence of
infrastructure that makes agent behavior predictable, verifiable, and reviewable.

% ============================================================================
\subsection{Standardized Conventions Eliminate Cognitive Overhead}

The system enforces comprehensive conventions for code organization,
naming, and structure that AI agents can internalize once and apply
universally. This standardization dramatically reduces the decision
space the AI must navigate when performing common development tasks.

Key conventions that benefit AI agents include:

\begin{itemize}
  \item \textbf{Predictable file organization}: Tests reside in \texttt{test/}
    directories adjacent to implementation code; golden files follow
    the pattern \texttt{test/outcomes/TestClass.test\_method/output/}

  \item \textbf{Systematic naming patterns}: Test classes follow
    \texttt{TestFooBar1} for testing class \texttt{FooBar}, test methods
    use \texttt{test\_calculate1()} for testing method \texttt{calculate()},
    with numbers enabling future expansion

  \item \textbf{Uniform test structure}: Every test adheres to the
    three-section pattern: prepare inputs, run test, check outputs,
    with mandatory docstrings and type annotations

  \item \textbf{Clear architectural boundaries}: Code is organized into
    well-defined components (e.g., DataPull, DataFlow, OMS, Market Data
    in the case of Causify system) with explicit dependency relationships and
    consistent internal structure
\end{itemize}

When asked to add tests for a new function \texttt{process\_data()} in
\texttt{foo/bar/module.py}, an AI agent immediately knows to create
\texttt{foo/bar/test/test\_module.py} containing class \texttt{Test\_process\_data1}
with appropriately structured test methods. No ambiguity, no need to
examine existing patterns, no opportunity for inconsistency. This
eliminates entire categories of trivial questions that would otherwise
require context switching or clarification.

% ============================================================================
\subsection{Container-Based Reproducibility Guarantees}

The Docker-based development system ensures that AI-generated code behaves
identically across all execution contexts, eliminating the entire class
of ``works on my machine'' failures that plague traditional development
workflows.

The containerization strategy provides several critical guarantees:

\begin{itemize}
  \item \textbf{Environment isolation}: All tests and builds execute
    in containers with precisely specified dependencies, eliminating
    host system variability as a source of failures

  \item \textbf{Version synchronization}: Container versions are locked
    and synchronized with code versions via \texttt{changelog.txt},
    ensuring that code and its execution environment remain compatible

  \item \textbf{Consistent tool behavior}: Dockerized executables work
    identically whether invoked from the host system, within development
    containers, or in CI/CD pipelines, using the sibling container
    pattern

  \item \textbf{Minimal host dependencies}: The thin environment approach
    requires only Docker and minimal Python packages on the host, enabling
    AI agents to operate in any environment (developer laptops, remote
    servers, CI/CD runners) with identical behavior
\end{itemize}

This reproducibility is particularly valuable for AI agents, which may
not have visibility into the host environment's configuration. When an
AI writes code that depends on specific Python packages or system
libraries, there is no ambiguity: the code either works correctly in the
container or it doesn't. The same container image used during development
is used in testing and production, eliminating environment-related
debugging cycles.

% ============================================================================
\subsection{Golden File Testing for Deterministic Verification}

The \texttt{hunitest.TestCase} framework's golden file testing pattern
via the \texttt{check\_string()} method provides AI agents with concrete,
deterministic feedback about code changes. Unlike traditional assertion-based
tests that provide binary pass/fail outcomes with limited context, golden
file tests expose the complete delta between expected and actual behavior.

This approach offers several advantages for AI-assisted development:

\begin{itemize}
  \item \textbf{Concrete, interpretable feedback}: When code changes
    affect behavior, the golden file diff shows exactly what changed,
    enabling the AI to assess whether the modification aligns with
    intent

  \item \textbf{Safe iteration cycles}: The AI can observe the precise
    impact of changes before committing to them, then use
    \texttt{pytest --update\_outcomes} to accept intentional modifications

  \item \textbf{Comprehensive regression detection}: Any unintended
    side effects from AI modifications appear immediately as unexpected
    diffs in golden files, even for distant parts of the system

  \item \textbf{Elimination of false positives}: Unlike complex boolean
    assertions that might pass due to implementation bugs, golden files
    capture complete output state, making incorrect passing tests far
    less likely
\end{itemize}

Consider an AI agent modifying a data processing pipeline: traditional
tests might only assert that certain summary statistics match expected
values, potentially missing subtle bugs. Golden file tests capture the
entire output, allowing both the AI and human reviewers to verify correctness
holistically rather than through narrow assertions.

% ============================================================================
\subsection{Layered Quality Gates for Incremental Validation}

The multi-tier quality enforcement system (pre-commit hooks, linters,
categorized tests, CI/CD pipelines) provides AI agents with progressive
validation checkpoints that enable efficient iteration while maintaining
quality standards.

The layered architecture delivers several benefits:

\begin{itemize}
  \item \textbf{Automated corrections}: Modifying linters (such as
    \texttt{Black}, \texttt{isort}, \texttt{autoflake}, \texttt{Prettier})
    automatically fix formatting and style issues, reducing cognitive
    burden for humans, allowing focus on logic and correctness

  \item \textbf{Immediate failure detection}: Pre-commit hooks catch
    syntax errors, file size violations, and accidentally committed
    secrets before humans and AI proceed to commit to the code repository

  % TODO(gp): Move this.
  \item \textbf{Graduated feedback cycles}: Fast tests (under 5 seconds
    per class) enable rapid iteration during initial development, while
    slow and superslow tests provide comprehensive validation before
    finalization

  \item \textbf{Specific, actionable diagnostics}: When quality gates fail,
    humans and AI receives precise error messages rather than generic failures
\end{itemize}

% ============================================================================
\subsection{Emergent AI-Optimization}

The Causify development system creates an ``AI-optimized codebase''
as an emergent property of designing for human developer productivity, team
scalability, and operational reliability. The system reduces the degrees of
freedom an AI agent must handle—eliminating questions like ``how should I
format this?'', ``where should tests live?'', ``how do I run them?''—while
providing rich, immediate feedback through multiple validation layers.

This architecture enables a more productive division of labor between
human developers and AI agents. The AI can focus on higher-level concerns
(algorithm correctness, business logic, architectural design) rather
than fighting with infrastructure, memorizing conventions, or debugging
environment-specific issues. The system's automation and quality gates
provide confidence that AI-generated code meets standards before human
review, while golden file tests and comprehensive test suites ensure
that subtle bugs or regressions are caught early.

Critically, these same properties that benefit AI agents also benefit
human developers, particularly those new to the codebase. The system's
discoverability, consistency, and automation reduce onboarding time
and cognitive load for all developers, whether human or artificial.
This suggests a broader principle: development infrastructure that
optimizes for clarity, automation, and reproducibility naturally becomes
more amenable to AI augmentation, creating a positive feedback loop where
better infrastructure enables more effective AI assistance, which in
turn accelerates infrastructure improvements.


  %- [1]
  %  [Mono vs. multi-repo](https://free.gitkraken.com/hubfs/Mono_v_Multi-Repo_debate_2023.pdf)
  %- [2]
  %  [Why Google stores billions of lines of code in a single repository](https://dl.acm.org/doi/10.1145/2854146)
  %- [3]
  %  [What it is like to work in Meta's (Facebook's) monorepo](https://blog.3d-logic.com/2024/09/02/what-it-is-like-to-work-in-metas-facebooks-monorepo/)
  %- [4]
  %  [Microsoft: How "Mono-repo" and "One Infra" Help Us Deliver a Better Developer Experience](https://devblogs.microsoft.com/appcenter/how-mono-repo-and-one-infra-help-us-deliver-a-better-developer-experience/)
  %- [5]
  %  [Uber: Faster Together: Uber Engineering's iOS Monorepo](https://www.uber.com/blog/ios-monorepo/)
\end{document}
